\documentclass[11pt,a4paper,openany]{book}
\usepackage[utf8]{inputenc}
\usepackage[T1]{fontenc}
\usepackage[margin=2.5cm]{geometry}
\usepackage{amsmath,amssymb,amsthm,mathtools}
\usepackage{booktabs,array,tabularx}
\usepackage{xcolor}
\usepackage{enumitem}
\usepackage{listings}
\usepackage{graphicx}
\usepackage{caption}
\usepackage[expansion=false]{microtype}
\usepackage{fancyhdr}
\usepackage{tikz}
\usetikzlibrary{arrows.meta,positioning,shapes.geometric,calc,fit,backgrounds}
\usepackage{pgfplots}
\pgfplotsset{compat=1.16}
\usepackage[most]{tcolorbox}
\usepackage[colorlinks=true,linkcolor=blue!50!black,urlcolor=blue!50!black,citecolor=blue!50!black]{hyperref}

\setlength{\parindent}{0pt}
\setlength{\parskip}{0.55em}
\setlist{itemsep=0.15em,topsep=0.3em}
\pagestyle{fancy}
\fancyhf{}
\fancyhead[LE]{\small\nouppercase{\leftmark}}
\fancyhead[RO]{\small\nouppercase{\rightmark}}
\fancyfoot[C]{\thepage}
\renewcommand{\headrulewidth}{0.3pt}

% ---------- notation ----------
\newcommand{\st}{\operatorname{stretch}}
\newcommand{\sq}{\operatorname{squash}}
\newcommand{\Pp}{\mathrm{P}}
\newcommand{\bpc}{\mathrm{bpc}}
\newcommand{\KL}{D_{\mathrm{KL}}}
\newcommand{\E}{\mathbb{E}}
\newcommand{\code}[1]{\texttt{#1}}
\DeclareMathOperator*{\argmin}{arg\,min}

% ---------- boxes ----------
\tcbset{enhanced jigsaw,breakable,boxrule=0.6pt,arc=2pt,left=6pt,right=6pt,top=4pt,bottom=4pt,fonttitle=\bfseries\sffamily}
\newtcolorbox[auto counter,number within=chapter]{example}[1][]{colback=blue!3,colframe=blue!55!black,title={Example \thetcbcounter\ifx&#1&\else\ --- #1\fi}}
\newtcolorbox{keyidea}[1][Key idea]{colback=green!4,colframe=green!45!black,title={#1}}
\newtcolorbox{pitfall}[1][Pitfall]{colback=red!4,colframe=red!55!black,title={#1}}
\newtcolorbox{handson}[1][Hands-on]{colback=orange!5,colframe=orange!60!black,title={#1}}
\newtcolorbox[auto counter,number within=chapter]{exercise}[1][]{colback=gray!5,colframe=gray!60!black,title={Exercise \thetcbcounter\ifx&#1&\else\ --- #1\fi}}
\theoremstyle{definition}
\newtheorem{definition}{Definition}[chapter]
\newtheorem{theorem}{Theorem}[chapter]
\newtheorem{proposition}[theorem]{Proposition}
\theoremstyle{remark}
\newtheorem*{remark}{Remark}

% ---------- code listings ----------
\definecolor{cmt}{rgb}{0.25,0.5,0.25}
\definecolor{kw}{rgb}{0.1,0.2,0.65}
\lstdefinestyle{ccode}{language=C,basicstyle=\ttfamily\footnotesize,keywordstyle=\color{kw}\bfseries,
  commentstyle=\color{cmt}\itshape,numbers=left,numberstyle=\tiny\color{gray},numbersep=6pt,
  breaklines=true,breakatwhitespace=false,showstringspaces=false,frame=single,framerule=0.3pt,
  rulecolor=\color{gray},xleftmargin=1.5em,tabsize=2,columns=fullflexible,keepspaces=true}
\lstdefinestyle{pseudo}{basicstyle=\ttfamily\footnotesize,numbers=none,frame=single,framerule=0.3pt,
  rulecolor=\color{gray},breaklines=true,columns=fullflexible,keepspaces=true,
  morekeywords={for,while,if,else,return,function,do,end,then,in}, keywordstyle=\color{kw}\bfseries}
\lstset{style=ccode}
\newcommand{\tikzbox}[3][fill=blue!8]{\node[draw,rounded corners=2pt,#1,minimum height=7mm,align=center,font=\small] (#2) {#3};}

\title{\Huge\bfseries Context Mixing from First Principles\\[0.4em]
\Large A Deep, Example-Driven Guide for Aspiring\\ Hutter Prize Competitors}
\author{}
\date{Compiled October 2026}

\begin{document}
\frontmatter
\maketitle

\chapter*{Preface}
\addcontentsline{toc}{chapter}{Preface}

This document is written for one purpose: to take you from ``I have heard that the best compressors use \emph{context mixing}'' to a state where you can \emph{design} a context-mixing compressor yourself, and understand every number it produces. It is deliberately long and deliberately concrete. Almost every idea is followed by a worked example with real numbers, and the last appendix contains a complete, working (about 200 lines of C) context-mixing compressor that I compiled and verified to round-trip, together with an ablation study showing what each component is worth.

\paragraph{How to read it.}
Part~I builds the mathematical foundations (information theory and arithmetic coding). Part~II builds a context-mixing compressor component by component: contexts, counters, bit histories, the logistic mixer, SSE/APM stages and the match model. Part~III moves to what is specific to the Hutter Prize: text and XML modelling for \code{enwik9}, dictionary preprocessing, neural components (the LSTM used in \code{cmix}), engineering under the time and memory limits, and a methodology for developing your own algorithm. The appendices contain the reference implementation, notation, hints and a reading list.

\paragraph{What is verified and what is not.}
Three kinds of statements appear in this document, and I try to keep them apart:
\begin{itemize}
\item \textbf{Mathematics} (derivations, worked examples). I recomputed every numeric example by program; they are exact up to the printed rounding.
\item \textbf{Experiments.} The tables in Appendix~A were produced by running the included code on a 1\,MB mixed English text sample from the machine where this document was prepared. This sample is \emph{not} Wikipedia, it contains a lot of near-duplicate license text, and so absolute numbers (especially the benefit of the match model) are \emph{much more favourable} than you will see on \code{enwik9}. Use them to learn the \emph{method} of ablation, not as predictions.
\item \textbf{Facts about the Prize and about other people's programs} (current record, rules, what \code{cmix}/\code{phda9}/STARLIT do). The record and rule figures were checked against the prize web pages in October 2026, but rules and records change; always re-read \url{http://prize.hutter1.net/}. Statements about the internals of specific third-party programs are from my background knowledge and the sources cited; where I am not sure I say so. \textbf{Before you copy any design decision, read the actual source code} (Appendix~E lists what to read).
\end{itemize}

\paragraph{Prerequisites.} Comfortable programming (C/C++ is the lingua franca of this field), basic probability, logarithms, and partial derivatives. Everything else is introduced.

\tableofcontents
\mainmatter

% =====================================================================
\part{Foundations}
% =====================================================================

\chapter{The Hutter Prize and the Big Picture}

\section{What the prize actually asks for}

The Hutter Prize (officially the ``500\,000\,\texteuro{} Prize for Compressing Human Knowledge'') rewards lossless compression of \code{enwik9}, the first $10^9$ bytes of an XML dump of the English Wikipedia taken in March 2006. The prize was announced in 2006; in February 2020 both the data (from the 100\,MB \code{enwik8} to the 1\,GB \code{enwik9}) and the prize pool (from 50\,000 to 500\,000 euros) were multiplied by ten.

\begin{keyidea}[The scoring rule]
You submit a compressor \code{comp9.exe} that turns \code{enwik9} into an archive \code{archive9.exe}, which is a \emph{self-extracting} program that reproduces \code{enwik9} exactly. Your score is the total
\[
S \;=\; \text{size}(\code{comp9.exe}\text{ as a zip}) \;+\; \text{size}(\code{archive9.exe}),
\]
and if $L$ is the previous record and $Z=500\,000$\texteuro{}, the award is $Z\,(L-S)/L$. The smallest claimable prize is 5\,000\texteuro, i.e.\ you must beat the previous record by at least $1\%$, and there is a 30-day public-comment period before payment.
\end{keyidea}

The important consequence of $S$ including the decompressor is that you cannot hide information in a huge pretrained model: every byte of dictionary, every learned parameter you ship counts against you, exactly like compressed data. This is the idea of the \emph{minimum description length} principle in action (Chapter~2).

\paragraph{Resource limits.} At the time of writing the official page lists, for recent records, about 10\,GB of RAM and about 50 hours of run time on a single core (and, per the rules mirrored in public README files, on the order of 100\,GB of disk). The exact limits and the exact reference hardware are set by the prize committee and ``may change without notice''; check the rules page before you design anything around them.

\section{The state of the art}

Table~\ref{tab:records} lists the \code{enwik9} records as published on the prize page and in public mirrors of it.

\begin{table}[h]
\centering
\small
\begin{tabular}{@{}llrrp{3.6cm}@{}}
\toprule
Author & Date & $S$ (bytes) & bits/byte & Program\\
\midrule
A.~Rhatushnyak & 2019 & 116\,673\,681 & 0.934 & phda9 v1.8 (pre-prize for enwik9)\\
A.~Margaritov & 31 May 2021 & 115\,352\,938 & 0.923 & STARLIT + cmix\\
S.~Kumar & 16 Jul 2023 & 114\,156\,155 & 0.913 & cmix-based\\
K.~Orav & 2 Feb 2024 & 112\,578\,322 & 0.900 & fx-cmix\\
K.~Orav \& B.~Knoll & 3 Sep 2024 & 110\,793\,128 & 0.886 & fx2-cmix\\
\bottomrule
\end{tabular}
\caption{Recent \code{enwik9} records (sizes include the decompressor). The last row is the record I could find as of October 2026; verify on the official page.}
\label{tab:records}
\end{table}

\begin{itemize}
\item \code{gzip -9} turns \code{enwik9} into roughly 323\,MB, i.e.\ $\approx 2.58$ bits per byte. The record is $110\,793\,128\cdot 8/10^9\approx 0.886$ bits per byte: about $2.9\times$ smaller than gzip.
\item Going from 116.67\,MB (2019) to 110.79\,MB (2024) is a total improvement of only $5.04\%$ in five years. Every further $1\%$ is $\approx 1.1$\,MB, and that is the minimum step to win money: to earn a prize one must reach $S\le 0.99\cdot 110\,793\,128\approx 109\,685\,196$ bytes.
\item This is a mature race. You will not win by a clever trick on top of gzip; you must understand and then \emph{extend} the state of the art, which today is a context-mixing ensemble with a neural network, tightly tuned for this exact file.
\end{itemize}

\section{The compute budget, quantified}

Because decompression must replay the model exactly, a context-mixing compressor spends the \emph{same} time to decode as to encode. So the 50-hour limit applies to the model itself:
\[
\frac{50\ \text{h}\cdot 3600\ \text{s/h}}{10^9\ \text{bytes}} = 180\ \mu\text{s per byte}\;=\;22.5\ \mu\text{s per bit}\;\approx\;5.6\ \text{KB/s}.
\]
For comparison a plain \code{lpaq}-style compressor processes on the order of a megabyte per second, so the prize allows more than two orders of magnitude more computation per byte than a fast CM. This is why the winning designs use hundreds of models and an LSTM: the budget is spent on \emph{more modelling}, not on entropy coding.

\section{Compression is prediction}

A deep fact, used throughout: \textbf{an adaptive model that assigns probability $p_t$ to the symbol that actually occurs at time $t$ can be turned, by arithmetic coding, into a compressor whose output length is $\sum_t -\log_2 p_t$ bits (plus a couple of bits).} The entropy coder is a solved, cheap, nearly lossless component. All the difficulty is in \emph{prediction}: the quality of the probabilities. This is why the prize is considered a test of ``intelligence'': a model that predicts Wikipedia text well must have absorbed syntax, semantics and facts. The prize FAQ states this motivation explicitly.

\section{A map of the compressor landscape}

\begin{description}
\item[Dictionary methods (LZ77, LZ78, LZMA).] Replace repeats with pointers. Fast, but only exploit exact repeats and code lengths are coarse.
\item[Symbol-ranking, BWT (bzip2).] Reorder the data to make it locally homogeneous, then code with simple adaptive models.
\item[PPM (Prediction by Partial Matching).] Predict with the longest context seen so far; if the symbol is new, ``escape'' to a shorter context. One context at a time is consulted.
\item[Context mixing (PAQ family, \code{cmix}).] Consult \emph{many} context models at once, each giving a probability for the next \emph{bit}; combine them with an adaptive, learned (online-trained) logistic function; refine with SSE; arithmetic-code the result. This is the subject of this document.
\item[Neural predictors (LSTM, Transformer).] A single learned network predicts the next symbol directly. Used as one more input to the mixer in \code{cmix}.
\end{description}

On the 100\,MB \code{enwik8} the first Hutter Prize baseline (Mahoney's \code{paq8f}, 2006) was 18\,324\,887 bytes, and the first prize-winning entry (Rhatushnyak's \code{paq8hp5}, September 2006) reached 17\,073\,018 bytes. The lineage \code{paq8f $\to$ paq8hp $\to$ phda9 $\to$ cmix $\to$ fx-cmix} is a straight line of context-mixing development.

\begin{figure}[h]
\centering
\begin{tikzpicture}[node distance=7mm and 6mm,>=Stealth,font=\small]
\tikzbox[fill=gray!15]{in}{input\\bits}
\tikzbox[right=of in]{ctx}{context\\hashing}
\tikzbox[right=of ctx]{cm}{$N$ context\\models}
\tikzbox[right=of cm,fill=green!12]{mix}{mixer\\network}
\tikzbox[right=of mix,fill=orange!15]{apm}{APM /\\SSE}
\tikzbox[right=of apm,fill=gray!15]{ac}{arithmetic\\coder}
\draw[->] (in)--(ctx); \draw[->] (ctx)--(cm); \draw[->] (cm)--node[above]{\scriptsize $p_i$}(mix);
\draw[->] (mix)--node[above]{\scriptsize $p$}(apm); \draw[->] (apm)--node[above]{\scriptsize $p'$}(ac);
\tikzbox[below=of mix,fill=purple!10]{lstm}{LSTM /\\neural models};
\tikzbox[below=of cm,fill=yellow!20]{match}{match\\model};
\draw[->] (lstm)--(mix); \draw[->] (match)--(mix);
\draw[->,dashed] (ac.south) |- ++(0,-1.6) -| node[pos=0.25,below]{\scriptsize coded bit $y$ updates everything} (ctx.south);
\end{tikzpicture}
\caption{The bit-level loop of a context-mixing compressor. After each bit is coded, \emph{every} component is updated with the true bit; the decoder runs the identical loop.}
\end{figure}

\section{Roadmap}

Chapters 2--3 give the theory of coding. Chapters 4--9 build a complete CM compressor. Chapters 10--12 specialise to Wikipedia text, neural components and engineering. Chapter 13 is a development methodology. If you are impatient, read Chapter~6 (mixing) and Appendix~A (the code) first, then return.


\chapter{Information Theory You Actually Need}

\section{Surprise, bits and code length}

If an event has probability $p$, observing it conveys $-\log_2 p$ bits of \emph{surprise}. A compressor that assigns probability $p$ to what actually happens is, in effect, ``paying'' $-\log_2 p$ bits for it.

\begin{table}[h]
\centering
\begin{tabular}{@{}rrr@{}}
\toprule
probability $p$ of what happened & cost $-\log_2 p$ (bits) & comment\\
\midrule
$1/2$ & 1.000 & fair coin\\
$0.9$ & 0.152 & expected event\\
$0.99$ & 0.0145 & near-certain\\
$0.999$ & 0.00144 & \\
$0.1$ & 3.322 & surprising\\
$0.01$ & 6.644 & very surprising\\
$0.001$ & 9.966 & \\
\bottomrule
\end{tabular}
\caption{Cost of an outcome as a function of the probability the model gave it.}
\end{table}

\begin{pitfall}[Overconfidence is expensive]
If the model says $P(\text{bit}=1)=0.999$ and the bit is $0$, the cost is $-\log_2 0.001=9.97$ bits. Each \emph{correct} prediction at that confidence costs only $0.0014$ bits, so a single miss wipes out the cost of about $6\,900$ correct predictions. The lesson, repeated throughout this document: \textbf{calibration} (stated probability $=$ actual frequency) matters as much as accuracy. SSE/APM stages (Chapter~7) exist precisely to fix calibration.
\end{pitfall}

\section{Entropy}

For a source with symbol probabilities $p(x)$ the \emph{entropy} is
\[
H(X)=-\sum_x p(x)\log_2 p(x)\quad\text{bits per symbol.}
\]
Shannon's source coding theorem: no lossless code can have average length below $H$, and codes exist that get within one bit (arithmetic coding gets within a fraction of a bit for the \emph{whole message}).

\begin{table}[h]
\centering
\begin{tabular}{@{}rr@{}}
\toprule
$P(1)$ & $H_b(P(1))$ bits\\
\midrule
0.5 & 1.0000\\
0.9 & 0.4690\\
0.99 & 0.0808\\
0.999 & 0.0114\\
\bottomrule
\end{tabular}
\caption{Binary entropy. A bit that is $99\%$ predictable carries only $0.08$ bits of information.}
\end{table}

\begin{example}[Entropy of ``abracadabra'']
The string \texttt{abracadabra} has 11 symbols: \texttt{a}$\times5$, \texttt{b}$\times2$, \texttt{r}$\times2$, \texttt{c}$\times1$, \texttt{d}$\times1$.
\begin{itemize}
\item \textbf{Order-0 model} (frequencies only): $H_0=-\tfrac{5}{11}\log_2\tfrac5{11}-2\cdot\tfrac2{11}\log_2\tfrac2{11}-2\cdot\tfrac1{11}\log_2\tfrac1{11}=2.0404$ bits/symbol, so $11\times2.0404=22.44$ bits, versus 88 bits as 8-bit ASCII.
\item \textbf{Order-1 model} (condition on the previous letter). The ten transitions are \texttt{ab}$\times2$, \texttt{br}$\times2$, \texttt{ra}$\times2$, \texttt{ac}, \texttt{ca}, \texttt{ad}, \texttt{da}. After \texttt{b}, \texttt{r}, \texttt{c} and \texttt{d} the next letter is \emph{always the same}, so those six transitions cost 0 bits. Only the context \texttt{a} is uncertain: it is followed by \texttt{b} twice, and by \texttt{c} and \texttt{d} once each, so the code lengths are $2\cdot1+1\cdot2+1\cdot2=6$ bits. The ideal order-1 cost of the ten transitions is therefore exactly $6.0$ bits, against $10\times2.04=20.4$ bits for order~0.
\end{itemize}
The lesson: richer context $\Rightarrow$ lower entropy \emph{if the statistics are known in advance}. A real compressor must \emph{learn} them while coding, which costs extra; that is the whole subject of adaptive modelling.
\end{example}

\section{Cross-entropy: the real cost of an imperfect model}

Let the data be generated by the true distribution $p$ but coded with a model $q$. The expected code length per symbol is the \emph{cross-entropy}
\[
H(p,q)=-\sum_x p(x)\log_2 q(x)=H(p)+\KL(p\,\|\,q),\qquad \KL(p\|q)=\sum_x p(x)\log_2\frac{p(x)}{q(x)}\ge0.
\]
The compressed size of a particular sequence $x_1\ldots x_n$ under a sequential model is exactly
\[
\boxed{\;\ell(x_{1:n})=\sum_{t=1}^n -\log_2 q(x_t\mid x_{1:t-1})\;}
\]
This is the \emph{only quantity you ever optimise}. Every technique in this document is a way of making this sum smaller on one specific file.

\begin{example}[Chain rule and sequential prediction]
Any joint distribution factorises as $q(x_{1:n})=\prod_t q(x_t\mid x_{<t})$. So building a compressor is equivalent to building a \emph{next-symbol predictor}. The arithmetic coder never needs to know anything else about the model.
\end{example}

\section{Mixtures of models and the cost of ``not knowing which is right''}

Suppose you have $K$ candidate models $q_1,\dots,q_K$ and no idea which is best on this file. A \emph{Bayesian mixture} $q_{\text{mix}}(x_{1:n})=\sum_k w_k q_k(x_{1:n})$ with prior weights $w_k$ (summing to 1) satisfies
\[
-\log_2 q_{\text{mix}}(x_{1:n}) \;\le\; -\log_2 q_k(x_{1:n}) + \log_2\frac1{w_k}\quad\text{for every }k.
\]
\emph{Proof.} $q_{\text{mix}}\ge w_kq_k$ since all terms are non-negative; take $-\log_2$. $\square$

So a mixture is never worse than the best model plus $\log_2(1/w_k)$ bits: with $K$ equal weights, a one-time cost of $\log_2K$ bits for ``choosing'' the best model, on the \emph{whole file}. This is an astonishingly good guarantee, and it is the theoretical reason to ensemble many models: with $K=1000$ models you pay only $\approx 10$ bits in total for not knowing which one is best.

\begin{keyidea}
Real context mixing improves on the Bayesian mixture in two ways: (i) the weights are learned \emph{continuously} and may depend on local context (the best model for XML tags differs from the best model for prose), and (ii) the combination is done in the logistic domain, which gives sharper predictions when models agree (Chapter~6). The price is that the simple $\log_2 K$ guarantee no longer applies; the benefit is that the compression is far better in practice.
\end{keyidea}

\section{Kolmogorov complexity, MDL and Hutter's motivation}

The Kolmogorov complexity $K(x)$ is the length of the shortest program that prints $x$. It is uncomputable, but any compressor gives an upper bound: $K(x)\le |\text{decompressor}|+|\text{archive}|$. The Hutter Prize's $S$ is exactly this upper bound for \code{enwik9}, with the decompressor included. Hutter's universal-AI theory (the \emph{AIXI} agent) is built on Solomonoff induction, which predicts with a mixture over \emph{all} computable models weighted by $2^{-\text{length}}$; context mixing is a practical, drastically restricted cousin of that idea. Whether or not you care about AGI, the take-away is practical: \textbf{model size is part of the score}, so a better model is only worth it if it saves more bytes than it costs to describe.

\section{How performance is measured}

\begin{itemize}
\item \textbf{bits per character} $\bpc = 8\,|\text{archive}|/|\text{original}|$. On text $\approx 1.6$--$2$ for good conventional compressors, $\approx 0.9$ for Hutter Prize records on \code{enwik9}.
\item \textbf{compression ratio} $=|\text{original}|/|\text{archive}|$ (the record has ratio about $9.03$).
\item For the Prize: the total $S$ in bytes, nothing else (subject to the resource limits).
\end{itemize}

\begin{exercise}[Cost of a streak]
A model predicts $P(1)=0.99$ for 1000 bits, and 10 of them turn out to be zeros. Compute the total code length in bits and compare with always predicting $0.5$. \emph{Answer:} $990\times0.0145+10\times6.644=14.4+66.4=80.8$ bits versus 1000 bits.
\end{exercise}

\begin{exercise}[Mixture bound]
Two models give a sequence probabilities $2^{-1000}$ and $2^{-1100}$. With equal prior weights, what does the mixture cost? \emph{Answer:} at most $1000+1=1001$ bits.
\end{exercise}

% =====================================================================
\chapter{Arithmetic Coding}
% =====================================================================

\section{The idea: nested intervals}

Represent the whole message as a number in $[0,1)$. Start with the interval $[0,1)$. For every bit, split the current interval in proportion to the model's probabilities and keep the sub-interval corresponding to the bit that actually occurred. After all bits, any number inside the final interval identifies the message; a number needs about $-\log_2(\text{width})$ bits to specify.

\textbf{Convention used here and in the code:} bit $1$ takes the \emph{lower} part of width $p_1\cdot(\text{width})$, bit $0$ takes the upper part.

\begin{example}[Encoding the bits 1\,0\,1\,1]
Model probabilities $P(1)=0.8,\,0.6,\,0.7,\,0.9$ for the four bits.
\begin{center}
\begin{tabular}{@{}cccccc@{}}
\toprule
step & $P(1)$ & bit $y$ & split point & new interval & width\\
\midrule
0 & & & & $[0,1)$ & 1\\
1 & 0.8 & 1 & $0+1\cdot0.8=0.8$ & $[0,\,0.8)$ & 0.8\\
2 & 0.6 & 0 & $0+0.8\cdot0.6=0.48$ & $[0.48,\,0.8)$ & 0.32\\
3 & 0.7 & 1 & $0.48+0.32\cdot0.7=0.704$ & $[0.48,\,0.704)$ & 0.224\\
4 & 0.9 & 1 & $0.48+0.224\cdot0.9=0.6816$ & $[0.48,\,0.6816)$ & 0.2016\\
\bottomrule
\end{tabular}
\end{center}
Final width $0.2016$, and indeed $-\log_2 0.2016=2.3104=0.3219+0.7370+0.5146+0.1520$, the sum of the four individual costs. Any number in $[0.48,0.6816)$ works; $0.5=0.1_2$ lies inside, so \emph{one bit} (``1'') suffices here (luck: the guaranteed bound is $\lceil-\log_2 w\rceil+1=4$ bits).

\textbf{Decoding the number $0.5$}, knowing the same probabilities:
\begin{center}
\begin{tabular}{@{}ccccc@{}}
\toprule
step & $P(1)$ & current value $v$ & test & decoded bit, rescaled $v$\\
\midrule
1 & 0.8 & 0.5 & $v<0.8$ & $1$; $v\leftarrow 0.5/0.8=0.625$\\
2 & 0.6 & 0.625 & $v\ge0.6$ & $0$; $v\leftarrow(0.625-0.6)/0.4=0.0625$\\
3 & 0.7 & 0.0625 & $v<0.7$ & $1$; $v\leftarrow0.0625/0.7=0.0893$\\
4 & 0.9 & 0.0893 & $v<0.9$ & $1$\\
\bottomrule
\end{tabular}
\end{center}
The decoder recovers $1011$. Note that \textbf{the decoder must compute exactly the same probabilities as the encoder}, bit for bit. This symmetry is why a CM decompressor costs as much time as the compressor, and why floating-point determinism matters (Chapter~11).
\end{example}

\begin{figure}[h]
\centering
\begin{tikzpicture}[x=11cm,y=1cm,font=\small]
\def\rows{{0,0.8,0.48,0.48}}
% interval bars: [lo,hi]
\foreach \lo/\hi/\yy/\lab in {0/1/0/{start}, 0/0.8/-1/{after 1}, 0.48/0.8/-2/{after 0}, 0.48/0.704/-3/{after 1}, 0.48/0.6816/-4/{after 1}}{
  \fill[blue!25] (\lo,\yy-0.28) rectangle (\hi,\yy+0.28);
  \draw (0,\yy-0.28) rectangle (1,\yy+0.28);
  \node[anchor=east] at (-0.01,\yy) {\lab};
}
\draw[red,thick] (0.5,0.5) -- (0.5,-4.5) node[below] {code point $0.5$};
\foreach \x in {0,0.25,0.5,0.75,1} \node[above] at (\x,0.35) {\scriptsize \x};
\end{tikzpicture}
\caption{The interval shrinks with every coded bit; its final width equals the product of the probabilities of the actual bits.}
\end{figure}

\section{Integer implementation (what you will actually write)}

Floating point is not used. The lpaq/paq8 coder keeps a 32-bit interval $[x_1,x_2]$ and a 12-bit probability $p\in[1,4095]$ that the next bit is 1:
\begin{lstlisting}[style=pseudo]
xmid = x1 + ((x2 - x1) >> 12) * p        // split point
if (y == 1) x2 = xmid else x1 = xmid + 1  // keep the sub-interval
while (((x1 ^ x2) & 0xff000000) == 0) {   // top bytes equal: they are final
    output_byte(x2 >> 24)
    x1 = x1 << 8;  x2 = (x2 << 8) | 255
}
\end{lstlisting}
The loop emits a byte as soon as the leading 8 bits of both ends agree, because no future bit can change them. At the end four bytes of $x_1$ are flushed. The decoder keeps the same $[x_1,x_2]$, reads the code bytes into $x$, and decodes $y=(x\le x_{mid})$.

\begin{example}[Integer coding with 12-bit probabilities]
Start with $x_1=\texttt{0x00000000}$, $x_2=\texttt{0xFFFFFFFF}$ and $p=3277$ ($\approx0.8\cdot4096$).
\begin{enumerate}
\item $(x_2-x_1)\gg12=\texttt{0xFFFFF}=1\,048\,575$; $x_{mid}=1\,048\,575\cdot3277=\texttt{0xCCCFF333}$. The bit is $y=1$, so $x_2\leftarrow\texttt{0xCCCFF333}$. The top bytes are $\texttt{00}\ne\texttt{CC}$: no output.
\item $p=2458$ ($\approx0.6\cdot4096$): $(x_2-x_1)\gg12=\texttt{0xCCCFF}$, $x_{mid}=\texttt{0x7AE84866}$. The bit is $y=0$, so $x_1\leftarrow x_{mid}+1=\texttt{0x7AE84867}$. Top bytes $\texttt{7A}\ne\texttt{CC}$.
\end{enumerate}
As real numbers: $\texttt{0x7AE84867}/2^{32}=0.4801$ and $\texttt{0xCCCFF333}/2^{32}=0.8000$, matching the interval $[0.48,0.8)$ of the previous example.
\end{example}

\section{Precision and its cost}

Two sources of inefficiency:
\begin{enumerate}
\item \textbf{Truncation} of $(x_2-x_1)\gg12$ and of $p$ to 12 bits. The interval is never less than $2^{24}$ wide after normalisation, so the relative error is tiny ($\lesssim2^{-12}$ in the worst case).
\item \textbf{Probability resolution.} With $p\le\frac{4095}{4096}$ the cheapest possible bit costs about $0.00035$ bits, i.e.\ $0.0028$ bits per byte. For ordinary text (0.9 bits per byte) that is irrelevant. For highly redundant data (e.g.\ long exact repeats) it is a measurable floor, which is why newer PAQ versions use 16-bit probabilities in the final stages. For \code{enwik9} the effect is small but non-zero; it is one of the many places where $0.1\%$ is hiding.
\end{enumerate}

\begin{pitfall}[Probability 0 and 1]
Never pass $p=0$ or $p=4096$ to the coder. A zero-probability event that occurs makes the interval empty and breaks decoding. Always clamp to $[1,4095]$ (or the equivalent in your precision).
\end{pitfall}

\section{Alternatives}

\begin{description}
\item[Range coders] are arithmetic coders working on bytes with multiplications on larger alphabets. Equivalent in efficiency.
\item[ANS (asymmetric numeral systems)] gives arithmetic-coding efficiency with table-driven speed. They are excellent for \emph{static} or slowly changing distributions (zstd, JPEG XL). They decode in reverse order of encoding, which is awkward when the probabilities depend on the data decoded so far, as in CM. For the Prize the coder is not the bottleneck; the binary arithmetic coder above is the standard choice.
\end{description}

\begin{exercise}
Encode the bits $0\,0\,1$ with $P(1)=0.3,0.2,0.5$ using the convention of this chapter and give the final interval and its width. \emph{Answer:} step 1 ($y=0$): split at $0.3$, interval $[0.3,1)$; step 2 ($y=0$): split at $0.3+0.7\cdot0.2=0.44$, interval $[0.44,1)$; step 3 ($y=1$): split at $0.44+0.56\cdot0.5=0.72$, interval $[0.44,0.72)$, width $0.28=0.7\cdot0.8\cdot0.5$.
\end{exercise}


% =====================================================================
\part{Building a Context-Mixing Compressor}
% =====================================================================

\chapter{Binary Decomposition and Contexts}

\section{Why predict one bit at a time?}

A byte has 256 possible values. One could predict a distribution over all 256, and PPM does exactly that. Context mixing instead decomposes each byte into 8 binary decisions, most significant bit first, and predicts each bit with a \emph{binary} model. Reasons:
\begin{enumerate}
\item Combining \emph{binary} predictions has a clean, learnable form (logistic mixing, Chapter~6). Combining 256-way distributions has no equally convenient form.
\item Each model needs only one number per context-and-bit-position, so memory and speed are predictable.
\item The binary arithmetic coder is simple and exact.
\end{enumerate}
The cost: 8 coding steps per byte instead of 1, and each model is consulted 8 times. At 180\,$\mu$s per byte (Chapter 1) this is affordable.

\section{The partial-byte context $c_0$}

When predicting bit number $k$ ($k=0\ldots7$) of a byte, the bits already coded for that byte are known. Prefix them with a leading 1 to get the number $c_0\in[1,255]$, which uniquely identifies the position in the binary tree of the byte.

\begin{example}[Coding the letter `e' ($=\texttt{0x65}=01100101_2$)]
\begin{center}
\begin{tabular}{@{}cccl@{}}
\toprule
bit index $k$ & value to code & $c_0$ \emph{before} coding & $c_0$ after\\
\midrule
0 & 0 & 1 ($=1_2$) & 2 ($=10_2$)\\
1 & 1 & 2 & 5 ($=101_2$)\\
2 & 1 & 5 & 11 ($=1011_2$)\\
3 & 0 & 11 & 22\\
4 & 0 & 22 & 44\\
5 & 1 & 44 & 89\\
6 & 0 & 89 & 178\\
7 & 1 & 178 & 357 ($=256+101$)\\
\bottomrule
\end{tabular}
\end{center}
After the 8th bit, $c_0=357$; subtracting 256 gives the byte value $101=\texttt{0x65}=\text{`e'}$. A context model keeps a \emph{separate} probability for each pair (context, $c_0$): up to 255 probabilities per context, but in practice only the few nodes on paths actually taken ever get touched.
\end{example}

\begin{figure}[h]
\centering
\begin{tikzpicture}[level distance=9mm,
  level 1/.style={sibling distance=48mm},level 2/.style={sibling distance=24mm},level 3/.style={sibling distance=12mm},
  every node/.style={font=\scriptsize,circle,draw,inner sep=1.5pt}]
\node {$c_0{=}1$}
 child {node {2} child {node {4} child{node{8}} child{node{9}}} child {node {5} child{node{10}} child{node{11}}}}
 child {node {3} child {node {6} child{node{12}} child{node{13}}} child {node {7} child{node{14}} child{node{15}}}};
\end{tikzpicture}
\caption{The first three levels of the binary tree over byte values. Left edge $=$ bit 0, right edge $=$ bit 1; $c_0$ of a child is $2c_0$ or $2c_0+1$.}
\end{figure}

\section{What is a context?}

A \emph{context} is any function of the past (and possibly of side information) that can be hashed to an index. The model remembers what happened \emph{every previous time} the same context value occurred. Good contexts partition the past into classes within which the next bit is similar.

\begin{example}[Contexts for predicting `a' in ``\ldots{}on the m'']
Suppose we have seen ``the cat sat on the mat. the cat sat on the m'' and are about to code the next byte (`a'):
\begin{center}
\begin{tabular}{@{}lll@{}}
\toprule
model & context value & why useful\\
\midrule
order 1 & \texttt{m} & after `m' vowels are likely\\
order 2 & \texttt{\textvisiblespace m} & word-initial `m'\\
order 3 & \texttt{e\textvisiblespace m} & ``(the) m''\\
order 4 & \texttt{he\textvisiblespace m} & more specific, seen once before\\
order 6 & \texttt{\textvisiblespace the m}\ & seen once before, followed by `a'\\
word & hash(``m'') + previous word ``the'' & word-level statistics\\
sparse & bytes 2 and 3 back, skipping 1 & robust to small differences\\
match & long match with the earlier ``on the m'' & predicts `a' very confidently\\
\bottomrule
\end{tabular}
\end{center}
Low orders give general, robust statistics; high orders give precise but rarely-available statistics. No single order is best. That is the motivation for mixing.
\end{example}

\paragraph{Families of contexts used in PAQ-style compressors.}
\begin{itemize}
\item \textbf{Order-$n$}: the last $n$ bytes ($n=1..6$ typical, plus a few high orders like 8, 12, 16, 24).
\item \textbf{Word contexts}: hash of the current partial word, of previous word(s), word bigrams; letter-case folded.
\item \textbf{Sparse / skip}: some of the last bytes, not contiguous. Useful for binary data and some text structure.
\item \textbf{Column / record}: position within a fixed-width record (tables, images).
\item \textbf{Indirect}: the context is the \emph{history of what followed} an earlier context (Section~\ref{sec:indirect}).
\item \textbf{Structure-aware (XML / markup)}: for \code{enwik9}: the enclosing tag, whether we are inside \verb|[[ ]]|, a table, a template (Chapter~10).
\end{itemize}

\section{PPM versus context mixing}

PPM picks one context, the longest one that has been seen, and uses an ``escape'' probability to fall back to shorter contexts when the symbol is new. It discards information from all other contexts. Context mixing consults \emph{all} orders simultaneously and lets the data decide, bit by bit, how much to trust each. On text, CM beats PPM by a large margin; in the early prize years the PPM-like \code{ppmd} variants were substantially behind \code{paq8}.

\section{Hashing and memory layout}

There are far more possible contexts than memory slots, so contexts are hashed into tables. Typical design (from lpaq/paq8):
\begin{itemize}
\item At each \emph{nibble} boundary (every 4 bits) compute a hash of (context hash, partial byte so far) and find the bucket of 16 or 64 bytes holding the states for the 15 or so tree nodes of the next nibble. This costs about \emph{two cache misses per byte per model}.
\item Each bucket carries a small checksum so that a collision (a different context mapping to the same bucket) is detected, rather than silently sharing statistics. On a mismatch the bucket is replaced (typically the lower-priority of two or more candidates, judged by the confidence of its first state).
\item Tables are sized as powers of two, and for \code{enwik9} the memory budget (about 10\,GB in total) is divided among models.
\end{itemize}

\begin{example}[Memory arithmetic]
A model on order-6 contexts over $10^9$ bytes can see up to $10^9$ distinct contexts. At 2 buckets per byte it makes $2\cdot10^9$ bucket visits. If you give it 1\,GB in 16-byte buckets you have only $6.7\cdot10^7$ buckets, so \emph{most} contexts must be evicted: the replacement policy decides what the model remembers. This is why careful tuning of table sizes and replacement rules gives real gains at \code{enwik9} scale, and why gains measured on \code{enwik8} may not carry over.
\end{example}

\begin{pitfall}[Hash collisions without checks]
The teaching compressor in Appendix~A omits checksums for clarity. With a $2^{22}$-slot table and a 1\,MB input collisions are rare; with $10^9$ bytes they would be devastating. Never use unchecked hashing at scale.
\end{pitfall}

% =====================================================================
\chapter{Predicting from One Context: Counters and Bit Histories}
% =====================================================================

\section{Plain frequency counts and the KT estimator}

The simplest predictor of the next bit in a context keeps two counts $n_0,n_1$. The \emph{Krichevsky--Trofimov (KT)} estimator is
\[
\hat p(1)=\frac{n_1+\tfrac12}{n_0+n_1+1}.
\]
It gives $0.5$ with no data, $0.75$ after one 1, $0.833$ after two 1's, $0.875$ after three. It never reaches 0 or 1, avoiding catastrophic costs. (The constant $\frac12$ is the optimal additive smoothing under a minimax criterion for a stationary Bernoulli source.)

\section{Adaptive-rate probability counters}

Storing two counts costs memory. A more compact alternative stores a probability $p$ and a small count $n$, updated by
\[
p\leftarrow p+\frac{y-p}{n+1.5},\qquad n\leftarrow\min(n+1,\;\text{limit}).
\]
This is the ``StateMap'' slot in Appendix~A (22-bit $p$ and 10-bit $n$ in one 32-bit word).

\begin{proposition}
With no limit on $n$ and $p_0=\frac12$, the update gives exactly $p_n=\dfrac{n_1+\frac14}{n+\frac12}$, where $n_1$ is the number of 1s among the first $n$ bits.
\end{proposition}
\emph{Proof by induction.} For $n=0$: $p_0=\frac{0.25}{0.5}=\frac12$. If $p_n(n+\frac12)=n_1+\frac14$ then $p_{n+1}=\dfrac{p_n(n+\frac12)+y}{n+\frac32}=\dfrac{n_1+y+\frac14}{(n+1)+\frac12}$. $\square$

So before the limit is reached the counter is a Dirichlet$(\tfrac14)$ estimator (slightly more aggressive than KT), and afterwards it becomes an exponential moving average with rate $\approx1/(\text{limit}+1.5)$. \textbf{The limit controls the stationarity assumption}: a high limit means ``this context is stationary, average everything''; a low limit means ``recent bits matter more''.

\begin{example}[Stationary versus adaptive]
A context sees eight 0s followed by four 1s. The probability that the \emph{next} bit is 1 after each bit:
\begin{center}
\small
\scriptsize\setlength{\tabcolsep}{3pt}
\begin{tabular}{@{}lcccccccccccc@{}}
\toprule
bit seen: & 0&0&0&0&0&0&0&0&1&1&1&1\\
\midrule
KT (counts, before update) & .500&.250&.167&.125&.100&.083&.071&.062&.056&.150&.227&.292\\
rate $1/(n{+}1.5)$, no limit & .167&.100&.071&.056&.045&.038&.033&.029&.132&.214&.283&.340\\
same, limit $n\le4$ & .167&.100&.071&.056&.045&.037&.030&.025&.202&.347&.466&.563\\
\bottomrule
\end{tabular}
\end{center}
(The first row gives the prediction made \emph{before} each bit is seen; the other two rows show $p$ \emph{after} the update.) After four 1s the stationary counter says $0.34$, while the fast-adapting counter says $0.56$: if the source really switched, the latter is closer to the truth; if the 1s were a random fluctuation the former is.
\end{example}

\section{Bit histories and non-stationary state machines}

Real data are \emph{non-stationary}: in a given context the distribution changes over time. Counts throw away \emph{order}. The sequences $000000\,11111$ and $010101\,01010$ have similar counts but very different implications.

PAQ replaces $(p,n)$ by a small \textbf{state} (one byte) which summarises the history: roughly ``$n_0$ zeros, $n_1$ ones, and the most recent bit'', with the constraint that counts are bounded and that when a bit is observed the \emph{opposite} count is reduced if it is large, so that the state forgets old evidence quickly when the data change. (The exact transition table in PAQ was generated by a program; there are about 255 states.) A simplified rule in the same spirit, illustrative only:
\begin{center}
observe $y$: $n_y\leftarrow n_y+1$; if $n_{1-y}>2$ then $n_{1-y}\leftarrow\lfloor n_{1-y}/2\rfloor+1$.
\end{center}
Starting from $(n_0,n_1)=(8,0)$ and observing four 1s the state moves $(5,1)\to(3,2)\to(2,3)\to(2,4)$. After three 1s it already says ``more 1s than 0s'', in agreement with the fast-adapting counter above and unlike the stationary counters.

\section{Indirect modelling: learning what a history means}
\label{sec:indirect}

The key insight of PAQ: do not convert the state to a probability with a fixed formula. Instead keep a \textbf{table mapping each state to a probability that is itself learned}, one table per model (an ``adaptive probability map'' or \emph{StateMap}):
\[
\text{context}\;\xrightarrow{\text{hash table}}\;\text{state}\;\xrightarrow{\text{StateMap}}\;p\;\xrightarrow{\text{stretch}}\;\text{mixer input}.
\]
The StateMap is trained on the \emph{outcomes}: whenever a bit $y$ is coded in a context whose state is $s$, the entry for $s$ is nudged toward $y$.

\begin{example}[What the map learns]
Consider the state ``three 1s, no 0s'' in an order-6 context of English text. KT would say $\hat p=0.875$. But in high-order contexts that have been completely deterministic so far, the true frequency of continued determinism is typically much higher (say $0.97$), while in an order-1 context the same state means little (perhaps $0.80$). Because each model has its \emph{own} StateMap, each learns the right number for \emph{its} situation, and the mixer weight handles the rest. (The numbers $0.97$ and $0.80$ are illustrative of the effect, not measurements.)
\end{example}

\begin{keyidea}[Why indirect beats direct]
A direct probability counter must learn each context's probability from scratch: after a single observation it has no idea whether a one-off event is a coincidence or a rule. An indirect model learns, once, across \emph{all} contexts of that model, what a one-off (or a streak) typically implies, and applies it immediately to new contexts. This is statistical borrowing of strength.
\end{keyidea}

\paragraph{The ``indirect context model''.} A different, also called \emph{indirect}, technique: for each order-$n$ context store the \emph{last few bytes that followed it} (e.g.\ the last two), then use that byte history as a context for a second-level model. It learns patterns like ``after this kind of pair the next byte is usually the same as the one before''.

\section{Beyond $\st(p)$: richer inputs per model}

PAQ8 gives the mixer more than one input per context:
\begin{itemize}
\item $\st(p)$: the log-odds (the main input).
\item $(p-\tfrac12)$ scaled: a \emph{linear-domain} input that gives the mixer a way to express different confidence shapes.
\item Special inputs for \emph{deterministic} contexts (only one bit value ever seen) and for \emph{run} information (the last byte seen in this context and how many times in a row).
\item A zero input for contexts that have never been seen: $\st(0.5)=0$, so an unseen context simply \emph{abstains} from the vote.
\end{itemize}

\begin{handson}[Mini-experiment]
Take the code in Appendix~A and replace each model's slot with a direct frequency counter (two 16-bit counts, KT estimate). Measure the size change on text. Then add a second input per model with $(p-0.5)$. This is exactly the kind of ablation that will tell you how much each idea is worth on \emph{your} data.
\end{handson}

\begin{exercise}[Counter limits]
Start from $p=0.5$, $n=0$ and apply the update $p\leftarrow p+(y-p)/(n+1.5)$ with limit 2 (so $n\le2$) to the bit sequence $1,1,1,1,0$. \emph{Answer:} $p=0.8333$ (rate $1/1.5$), $0.9$ (rate $1/2.5$), $0.9286$ (rate $1/3.5$); from here $n$ stays frozen at 2, so the rate stays $1/3.5=0.2857$: $p=0.9490$ after the fourth 1, and $p=0.9490-0.2857\cdot0.9490=0.6779$ after the final 0.
\end{exercise}


% =====================================================================
\chapter{The Heart of Context Mixing: Logistic Mixing}
% =====================================================================

\section{The problem}

At each bit we have $N$ models, each giving a probability $p_i=P_i(y=1)$. We need \emph{one} probability $p$. We want a rule that (a) is cheap, (b) can express ``trust model 3 a lot when the context is XML'', (c) can be \emph{trained online} as the data stream by, and (d) combines evidence sensibly.

\section{Stretch and squash}

\begin{definition}
The \emph{logit} or \emph{stretch} function and its inverse, the \emph{logistic} or \emph{squash} function:
\[
\st(p)=\ln\frac{p}{1-p},\qquad \sq(x)=\frac1{1+e^{-x}},\qquad \sq(\st(p))=p.
\]
\end{definition}
$\st(p)$ is the log-odds of the event. It maps probabilities $(0,1)$ to the whole real line, with $\st(\tfrac12)=0$, positive when $1$ is more likely, negative when $0$ is, and symmetric: $\st(1-p)=-\st(p)$.

\begin{table}[h]
\centering
\begin{tabular}{@{}rrr@{}}
\toprule
$p$ & $\st(p)$ & fixed-point value $\lfloor 256\,\st(p)\rceil$ (as in the code)\\
\midrule
0.001 & $-6.907$ & $-1768$\\
0.01 & $-4.595$ & $-1176$\\
0.1 & $-2.197$ & $-562$\\
0.25 & $-1.099$ & $-281$\\
0.5 & $0$ & $0$\\
0.75 & $1.099$ & $281$\\
0.9 & $2.197$ & $562$\\
0.99 & $4.595$ & $1176$\\
0.999 & $6.907$ & $1768$\\
\bottomrule
\end{tabular}
\caption{Stretch values. The integer code clamps the domain to $\pm2047$ ($=\pm8$ in real units), i.e.\ $p\in[0.00034,\,0.99966]$.}
\end{table}

\begin{figure}[h]
\centering
\begin{tikzpicture}
\begin{axis}[width=0.47\textwidth,height=5cm,xlabel={$x$},ylabel={$\sq(x)$},domain=-8:8,grid=major,title={squash}]
\addplot[blue,thick,samples=60]{1/(1+exp(-x))};
\end{axis}
\end{tikzpicture}\hfill
\begin{tikzpicture}
\begin{axis}[width=0.47\textwidth,height=5cm,xlabel={$p$},ylabel={$\st(p)$},domain=0.002:0.998,ymin=-7,ymax=7,grid=major,title={stretch}]
\addplot[red,thick,samples=100]{ln(x/(1-x))};
\end{axis}
\end{tikzpicture}
\caption{Squash and stretch. Stretch magnifies the region near 0 and 1 where the \emph{costly} mistakes live.}
\end{figure}

\section{Logistic mixing}

\begin{definition}[Logistic mixing]
Given inputs $s_i=\st(p_i)$ and weights $w_i$ (one weight vector per \emph{weight set}, chosen by a small context), the mixed probability is
\[
\boxed{\;p=\sq\Big(\sum_{i=1}^Nw_is_i\Big)\;}
\]
\end{definition}
This is a single artificial neuron with a sigmoid output (logistic regression). Its training signal is the very quantity we want to minimise, the coding cost.

\subsection{Why in the logistic domain? The product-of-experts view}

Define the \emph{geometric} mixture of the experts' distributions: $P(y)\propto\prod_iP_i(y)^{w_i}$. For a binary outcome,
\[
\ln\frac{P(1)}{P(0)}=\sum_iw_i\ln\frac{P_i(1)}{P_i(0)}=\sum_iw_is_i,
\]
so $P(1)=\sq(\sum w_is_i)$. Logistic mixing is exactly the (normalised) weighted geometric mean of the models' distributions, a \emph{product of experts}. Contrast with the \emph{linear} mixture $p=\sum_i\lambda_ip_i$.

\begin{example}[Two confident, independent models]
Two models each predict $P(1)=0.99$.
\begin{itemize}
\item \textbf{Linear mixing} with $\lambda=(\frac12,\frac12)$: $p=0.99$. Two agreeing experts give no extra confidence.
\item \textbf{Logistic mixing} with $w=(1,1)$: $\sq(2\cdot4.595)=\sq(9.19)=0.9999$. Cost if the bit is 1: $0.00014$ bits instead of $0.0145$.
\end{itemize}
If the two models see \emph{independent} evidence (different contexts that both say ``1''), adding log-odds is Bayes' rule (naive Bayes): the evidence accumulates. The learned weights $w_i$ then correct for the fact that real models are \emph{not} independent (weights below 1) or that some models are more reliable (weights above 1). Linear mixing can never produce a probability more extreme than its most extreme input, which is a severe limitation for highly predictable data like text.
\end{example}

\begin{example}[Disagreement]
Model A says $0.9$, model B says $0.1$ with equal weights $w=(0.5,0.5)$. Then $s=(2.197,-2.197)$ and $p=\sq(0)=0.5$: a sensible ``don't know''. With $w=(0.8,0.3)$ the sum is $0.8\cdot2.197-0.3\cdot2.197=1.10$ and $p=0.75$: trust A more. Weights may also be \emph{negative}: a model that is systematically wrong in some context gets a negative weight and its vote is inverted. Linear mixing with non-negative weights cannot do this.
\end{example}

\section{Learning the weights online}

Let the loss be the coding cost of the bit, $\ell=-\ln P(y)=-y\ln p-(1-y)\ln(1-p)$ (natural log; the base only scales the learning rate), with $p=\sq(z)$, $z=\sum_iw_is_i$.

\begin{theorem}[Gradient of the coding cost]
$\displaystyle\frac{\partial\ell}{\partial w_i}=-(y-p)\,s_i.$
\end{theorem}
\emph{Proof.} $\frac{d\sq}{dz}=p(1-p)$, and $\frac{\partial\ell}{\partial p}=-\frac yp+\frac{1-y}{1-p}=\frac{p-y}{p(1-p)}$. By the chain rule $\frac{\partial\ell}{\partial z}=p-y$, and $\frac{\partial z}{\partial w_i}=s_i$. $\square$

So \textbf{online gradient descent} is simply
\[
\boxed{\;w_i\leftarrow w_i+\eta\,(y-p)\,s_i\;}
\]
where $\eta$ is the learning rate. Read it aloud: ``increase the weight of every model that voted in the direction of the actual bit, in proportion to how strongly it voted and how wrong the mixture was.''

\begin{keyidea}[Two lucky facts]
(1) Because the loss is the log-loss and the output is a sigmoid, the awkward factor $p(1-p)$ cancels. The gradient does not vanish when the mixer is confidently wrong. (With squared-error loss it would; this is one reason PAQ is trained with log-loss.) (2) The loss is \emph{convex} in $w$, so there are no local minima and online gradient descent enjoys regret guarantees of order $\sqrt T$ against the best fixed weight vector in hindsight.
\end{keyidea}

\section{A complete numerical example}

\begin{example}[One mixing step, by hand]
Three models predict $P(1)=0.9,\,0.6,\,0.3$. Current weights $w=(0.4,0.3,0.3)$, learning rate $\eta=0.02$.

\textbf{Forward.} $s=(\st(0.9),\st(0.6),\st(0.3))=(2.1972,\,0.4055,\,-0.8473)$. The dot product is
\[
\begin{aligned}z&=0.4\cdot2.1972+0.3\cdot0.4055+0.3\cdot(-0.8473)\\&=0.8789+0.1216-0.2542=0.7463,\end{aligned}\qquad p=\sq(0.7463)=0.6784.
\]
(A plain weighted average of the probabilities would have given $0.4\cdot0.9+0.3\cdot0.6+0.3\cdot0.3=0.63$.)

\textbf{Case A: the bit is $y=1$.} Cost before: $-\log_2 0.6784=0.5598$ bits. Error $y-p=0.3216$. Weight updates $\Delta w_i=0.02\cdot0.3216\cdot s_i$:
\[
\Delta w=(+0.0141,\;+0.0026,\;-0.0055),\qquad w'=(0.4141,\,0.3026,\,0.2945).
\]
Model 1 (which strongly said ``1'') gains weight; model 3 (which said ``0'') loses it. Recomputing with the new weights gives $p'=0.6863$, cost $0.5430$ bits: after \emph{one} update the same inputs would have been coded $0.017$ bits cheaper.

\textbf{Case B: the bit is $y=0$.} Cost before: $-\log_2(1-0.6784)=1.6366$ bits. Error $y-p=-0.6784$, $\Delta w=(-0.0298,\,-0.0055,\,+0.0115)$, $w'=(0.3702,\,0.2945,\,0.3115)$, new $p'=0.6612$, cost $1.5617$ bits. Model 1 is penalised for having been confidently wrong.
\end{example}

\section{Fixed-point implementation}

The compressor in Appendix~A is integer-only (deterministic and fast):
\begin{itemize}
\item inputs $s_i=\text{stretch}(p_i)\in[-2047,2047]$ (real value $\times256$);
\item weights $w_i$ in 16.16 fixed point ($65536=1.0$), initialised to $2^{14}=0.25$;
\item output $p=\text{squash}\big((\sum_is_iw_i)\gg16\big)$, a 12-bit probability;
\item \code{err = ((y<<12) - p) * lr}, then \code{w[i] += (s[i]*err + 0x8000) >> 16}. 
\end{itemize}
Working through the scalings (inputs $\times256$, weights $\times2^{16}$, error $\times4096\cdot lr$, final shift $2^{-16}$) the integer update equals the floating-point rule $w_i\leftarrow w_i+\eta(y-p)s_i$ with
\[
\eta=\frac{lr}{4096}.
\]
So the value $lr=6$ used in Appendix~A corresponds to $\eta\approx0.0015$; practical integer rates lie roughly in $2$--$16$, i.e.\ $\eta\approx0.0005$--$0.004$. (The learning rate $0.02$ in the hand example above was chosen large so that the weight changes are visible at four decimals.)

\begin{pitfall}[The learning rate is the most sensitive parameter]
Too small: the mixer cannot follow changes in the data. Too large: the weights jitter and the code length rises (the ablation in Appendix~A shows that disabling learning altogether costs 17\%). Good compressors use a \emph{decaying} or \emph{adaptive} rate, different per mixer, tuned by experiment; some modern variants replace plain SGD by RMS-normalised or Adam-style updates.
\end{pitfall}

\section{Weight sets selected by context}

A single weight vector must be right everywhere. Instead keep \emph{many} weight vectors and use a small context $m$ to choose which one to use and train:
\[
p=\sq\Big(\sum_iw_{m,i}s_i\Big).
\]
The selector $m$ can be the partial byte $c_0$ (the best model mix differs between the first bit of a byte, which is nearly deterministic for ASCII text, and the last), the length of the current match, the previous byte, the bit position, the current XML state\ldots{} The trade-off is dilution: with $M$ sets each trains on $1/M$ of the data. Typical designs use several mixers with \emph{different} selectors and combine them (next section).

\section{Two-layer mixing networks}

PAQ8 uses a small neural network of two layers:
\begin{enumerate}
\item \textbf{Layer 1:} $K$ mixers, each using the same $N$ inputs but a different selector context $m_k$. Each outputs $p^{(k)}=\sq(\sum_iw_{m_k,i}s_i)$.
\item \textbf{Layer 2:} a final mixer takes $\st(p^{(k)})$, $k=1..K$ as inputs, with its own weight set (selected by another small context), and produces the final probability.
\end{enumerate}
Each layer-1 mixer is trained to minimise the coding cost \emph{on its own} output; the final mixer is trained on the final output.

\begin{figure}[h]
\centering
\begin{tikzpicture}[>=Stealth,font=\small,inp/.style={draw,circle,inner sep=1.5pt,fill=blue!8},mx/.style={draw,rounded corners,fill=green!12,minimum width=14mm,minimum height=6mm}]
\foreach \i in {1,...,5} { \node[inp] (i\i) at (0,-\i*0.8+3.2) {\scriptsize $s_\i$}; }
\node[mx] (m1) at (3.5,2) {mixer 1};
\node[mx] (m2) at (3.5,0.6) {mixer 2};
\node[mx] (m3) at (3.5,-0.8) {mixer 3};
\node[mx,fill=orange!20] (f) at (6.8,0.6) {final};
\foreach \i in {1,...,5} \foreach \m in {1,2,3} \draw[->,gray!60] (i\i)--(m\m);
\foreach \m in {1,2,3} \draw[->] (m\m)--(f);
\draw[->] (f)--++(1.6,0) node[right] {$p$};
\node[align=center,font=\scriptsize] at (3.5,-1.9) {selected by $c_0$, match length,\\order-1 byte, \ldots};
\end{tikzpicture}
\caption{A two-layer mixer: several mixers with different weight-set selectors, then a small final mixer.}
\end{figure}

Why does it help? Each first-layer mixer is a \emph{specialist} for one way of slicing the data; the second layer learns which specialist to trust at the moment. It is an ensemble of ensembles. In practice adding a second layer and a handful of well-chosen selectors gives a few percent.

\section{Ablation results from the reference implementation}

I ran the compressor of Appendix~A, switching off one component at a time, on a 1\,MB English text sample (strongly repetitive license/documentation text; see the Preface for the caveat).

\begin{table}[h]
\centering
\begin{tabular}{@{}lrr@{}}
\toprule
configuration & size (bytes) & vs.\ full\\
\midrule
full model (orders 0,1,2,3,4,6 + word + match + 2 APMs) & 89\,778 & --\\
no word model & 91\,847 & $+2.3\%$\\
single weight set (instead of 256 selected by $c_0$) & 90\,404 & $+0.7\%$\\
no APM stages & 90\,925 & $+1.3\%$\\
no order-6 context & 90\,190 & $+0.5\%$\\
no orders 3, 4, 6 & 95\,424 & $+6.3\%$\\
mixer learning disabled (fixed weights 0.25) & 104\,626 & $+16.5\%$\\
no match model & 133\,226 & $+48\%$\\
\bottomrule
\end{tabular}
\caption{Ablation on a 1\,MB sample. The enormous value of the match model is a property of \emph{this} sample (many near-duplicate documents); on Wikipedia it is far smaller, though still essential.}
\end{table}

\begin{center}
\begin{tikzpicture}
\begin{axis}[xbar,width=0.9\textwidth,height=6.3cm,xlabel={size increase over full model (\%)},symbolic y coords={no match,learning off,no orders 3/4/6,no word,no APM,single weight set,no order 6},ytick=data,y dir=reverse,nodes near coords,every node near coord/.append style={font=\scriptsize},xmin=0,xmax=58,bar width=7pt]
\addplot coordinates {(48.4,no match) (16.5,learning off) (6.3,no orders 3/4/6) (2.3,no word) (1.3,no APM) (0.7,single weight set) (0.5,no order 6)};
\end{axis}
\end{tikzpicture}
\end{center}

\begin{exercise}[Reproduce the example]
Compute by hand the mixed probability for inputs $p=(0.8,0.7)$ with weights $w=(0.6,0.6)$. \emph{Answer:} $\st=(1.3863,0.8473)$; $z=0.6\cdot2.2336=1.3402$; $p=0.7926$. Note this is \emph{higher} than either input would suggest under averaging (0.75), because the two models agree.
\end{exercise}

\begin{exercise}[Gradient check]
In the example above with $y=1$ show that the cost decreases for a small step in the negative gradient direction. \emph{Hint:} the first-order change in cost is $-\eta\|(y-p)s\|^2<0$.
\end{exercise}


% =====================================================================
\chapter{SSE / APM: Fixing Calibration After Mixing}
% =====================================================================

\section{Idea}

The mixer's output can be systematically miscalibrated: perhaps, in some situation, when it says $0.62$ the bit is $1$ only $55\%$ of the time, or $75\%$. A \textbf{secondary symbol estimation} (SSE) or \textbf{adaptive probability map} (APM) stage learns a \emph{correction function} $p\mapsto p'$ \emph{for each small context}. It is a learned calibration curve.

Implementation (as in Appendix~A): for each context $m$ keep $33$ values $t_{m,0..32}$, the corrected probabilities at the 33 stretch values $-8,-7.875,\ldots,+8$ (spacing 128 in fixed-point). To refine a probability $p$:
\begin{enumerate}
\item compute $s=\st(p)$ in fixed-point, locate it between two neighbouring nodes $j,j+1$ with interpolation weight $\omega=(s-s_j)/128$;
\item output $p'=(1-\omega)\,t_{m,j}+\omega\,t_{m,j+1}$;
\item after coding bit $y$, move the \emph{nearer} node toward $y$: $t\leftarrow t+(y\text{-target}-t)/2^{r}$ (rate $r=7$).
\end{enumerate}
Nodes are initialised to the \emph{identity} map $t_{m,j}=\sq((j-16)\cdot128)$, so at the beginning the APM changes nothing.

\begin{example}[One APM lookup and update]
The mixer outputs $p=0.62$. In fixed-point $s=\lfloor256\st(0.62)\rceil=125$. Then $(s+2048)\cdot32=69\,536$; the node index is $69\,536\gg12=16$ and the interpolation weight is $69\,536-16\cdot4096=4000$ out of $4096$ ($=0.977$), i.e.\ the input lies $97.7\%$ of the way from node 16 ($s=0$) to node 17 ($s=128$). With the identity initialisation $t_{16}=2048\cdot16$ and $t_{17}=2549\cdot16$ (units: 12-bit probability $\times16$):
\[
p'=\frac{32768\cdot96+40784\cdot4000}{2^{16}}=2537\ \Rightarrow\ \frac{2537}{4096}=0.619.
\]
Now the bit is $y=1$. The nearer node is $17$ (index $16+(4000\gg11)=17$). The update target is $g=2^{16}+2^{7}-2=65662$ and
$t_{17}\leftarrow40784+\big((65662-40784)\gg7\big)=40784+194=40978$.
The same input now yields $p'=2549/4096=0.6223$. If, in this context, mixer outputs near $0.62$ really are followed by 1 about 75\% of the time, the node values will drift upward over time and the APM will output $\approx0.75$.
\end{example}

\section{Chaining and averaging}

PAQ-style final stages typically use two or more APMs with \emph{different} contexts (e.g.\ $c_0$; the previous byte and $c_0$; a hash of the order-2 context; the match state) and then average the outputs with the original, e.g.\ $p_{final}=\tfrac14(2p_{mix}+p_1+p_2)$ as in Appendix~A. The average is a safety measure: an APM with a sparse context is noisy and averaging reduces its variance. In the ablation, removing both APMs costs $1.3\%$.

\begin{keyidea}
The mixer is trained to be good \emph{on average} over its weight-set selector; the APM corrects the residual calibration error \emph{conditionally on additional context}, including contexts the mixer never saw. SSE therefore adds information, not just calibration.
\end{keyidea}

% =====================================================================
\chapter{The Match Model and Long-Range Structure}
% =====================================================================

\section{Why a separate model for long repeats}

Context models of order up to about 6 are dense. A context of order 20 or 32 is \emph{almost always novel} unless a long string repeats, and if it repeats, the next symbol is very likely to repeat too. The \textbf{match model} handles this cheaply: instead of a hash table of statistics per context, it keeps one pointer into the history.

\section{Algorithm}

\begin{lstlisting}[style=pseudo]
after each byte c is coded and appended to buf[0..pos):
    if (len > 0 and buf[ptr] == c)  { len++;  ptr++ }       // keep following
    else len = 0
    h = hash(last MINLEN bytes)
    if (len == 0) {
        ptr = table[h]                                       // most recent occurrence
        if (ptr > 0) while (len < MAXLEN and buf[ptr-len-1] == buf[pos-len-1]) len++
    }
    table[h] = pos

for each bit of the next byte:
    if (len > 0 and the bits coded so far agree with buf[ptr]):
        expected = next bit of buf[ptr]
        input = stretch( StateMap[min(len,31)][expected] )   // learned confidence
    else input = 0, len = 0 for this byte
\end{lstlisting}

Two details deserve emphasis.

\paragraph{The confidence is learned, not hard-coded.} The StateMap indexed by (match length bucket, expected bit) learns how reliable a match of a given length is. In my test the learned reliabilities, measured as the fraction of predicted bits that were correct, were:

\begin{center}
\begin{tikzpicture}
\begin{axis}[width=0.85\textwidth,height=5.5cm,xlabel={match length (bytes, capped at 31)},ylabel={fraction of predicted bits correct},ymin=0.94,ymax=1.0,grid=major,mark size=1.5pt]
\addplot[blue,mark=*] coordinates {(6,0.9502)(7,0.9598)(8,0.9664)(9,0.9677)(10,0.9704)(11,0.9763)(12,0.9806)(14,0.9841)(16,0.9832)(18,0.9821)(20,0.9859)(22,0.9889)(24,0.9898)(26,0.9888)(28,0.9922)(30,0.9888)(31,0.9986)};
\end{axis}
\end{tikzpicture}
\end{center}

The curve rises smoothly (a match of 6 bytes is right 95\% of the time, but a match of 31 or more is right 99.86\% of the time). \emph{This is a property of my highly repetitive test file}; on Wikipedia the curve is lower and flatter, and you should measure it yourself. The point is that the mixer receives $\st(0.9986)=6.6$ when the match is long and a much smaller number for short matches, so the mixer is effectively told how much to trust the match.

\paragraph{A mismatch cancels the match for the rest of the byte.} Once a predicted bit is wrong, the pointer is dropped (and re-found at the next byte boundary through the hash). Production implementations add \emph{recovery}: they remember the old pointer for a few bytes in case the text diverges briefly (a typo, a changed number) and continues the same.

\begin{example}[A match in action]
History: \texttt{the quick brown fox jumps over the lazy dog. the quick brown\textvisiblespace}, and suppose the text then continues with an `x' (``\ldots{}brown x\ldots{}'') instead of the `f' of the first occurrence. After the second ``\texttt{brown\textvisiblespace}'' the hash of the last six bytes (\texttt{brown\textvisiblespace}) finds the earlier occurrence, and the backward comparison extends the match to the 16 bytes ``\texttt{the quick brown\textvisiblespace}'' (it stops at the beginning of the history). The predicted next byte is `f' ($=\texttt{0x66}=01100110_2$), the actual byte is `x' ($=\texttt{0x78}=01111000_2$). Bits 0--2 ($0,1,1$) agree with the prediction and the match model votes confidently (reliability for length 16). At bit 3 the prediction (0) is wrong; the cost is the usual large penalty for a confident miss, the pointer is dropped for the rest of the byte, and the match is re-acquired through the hash after the byte ends.
\end{example}

\section{Making the match model smarter}

\begin{itemize}
\item \textbf{Several match models} with different minimum lengths (e.g.\ 4, 8, 20, 32), each with its own pointer: the short one finds recent-but-weak matches, the long one finds reliable ones.
\item \textbf{Verify the candidate}: store a few candidate positions per hash and choose the longest actual match.
\item \textbf{Expected byte as a context}: use the predicted byte (and the match length bucket) as the context of an ordinary context model and as the selector for a mixer weight set. That lets the system learn: ``when the match predicts a vowel it is right, when it predicts a digit it often fails''.
\item \textbf{Use it in the APM}: SSE indexed by (expected bit, length bucket).
\end{itemize}

\section{Long-range structure in Wikipedia}

The XML dump contains massive structural repetition: the per-article page skeleton (page, title, id and revision tags), contributor lines, boilerplate in templates, repeated link text, tables and lists in similar formats. A match model sees this immediately, while order-$n$ models for $n\le6$ see only fragments. This is also where \emph{article ordering} matters, because the match model prefers related text to be close: see Chapter~10.

% =====================================================================
\chapter{Putting It Together: One Bit, Start to Finish}
% =====================================================================

\section{The full loop}

\begin{lstlisting}[style=pseudo]
for each byte:
    for each bit position:
        for each model i:   p_i = lookup(context_i, c0)            // state -> StateMap -> prob
                            s_i = stretch(p_i) (plus extra inputs)
        for each mixer k:   z_k = sum_i w[k][sel_k][i] * s_i ; p_k = squash(z_k)
        p_mix = final_mixer(stretch(p_k))
        p = average( p_mix, APM_1(p_mix, ctx_1), APM_2(p_mix, ctx_2) )
        encode(bit, p)                                             // or decode
        --- both sides now know the bit y ---
        for each model:     update state/counter for the bit
                            update the StateMap entry toward y
        for each mixer:     w += eta * (y - p_k) * s_i
        update APMs, c0 = 2*c0 + y
    at byte boundary: update histories, hashes, match model, word hash, ...
\end{lstlisting}

\section{A numerical walk-through of a single bit}

\begin{example}[Predicting the third bit of `e' after ``th'']
We are coding `e' ($01100101_2$) in the context ``\ldots{}in th''. We are at bit index 2 with $c_0=5$ (bits $01$ already coded). The next bit is $y=1$. Models report $P(y=1)$:
\begin{center}
\begin{tabular}{@{}llrrrr@{}}
\toprule
input & source & $p_i$ & $s_i=\st(p_i)$ & weight $w_i$ (set $c_0{=}5$) & cost if $y{=}1$ (bits)\\
\midrule
1 & order 0 & 0.55 & 0.2007 & 0.10 & 0.862\\
2 & order 1 (`h') & 0.80 & 1.3863 & 0.20 & 0.322\\
3 & order 2 (``th'') & 0.93 & 2.5867 & 0.30 & 0.105\\
4 & order 3 (`\textvisiblespace th') & 0.97 & 3.4761 & 0.25 & 0.044\\
5 & word model & 0.90 & 2.1972 & 0.15 & 0.152\\
6 & match (inactive) & -- & 0 & 0.00 & --\\
7 & bias (constant 1) & -- & 1 & $-0.05$ & --\\
\bottomrule
\end{tabular}
\end{center}
\textbf{Mixer:} $z=0.10\cdot0.2007+0.20\cdot1.3863+0.30\cdot2.5867+0.25\cdot3.4761+0.15\cdot2.1972-0.05\cdot1=2.2219$, so $p_{mix}=\sq(2.2219)=0.9022$. (The best single model alone would cost $0.044$ bits, the mixture costs $0.1485$ bits \emph{because the weights are not yet trained to trust order 3 more}; this is just a snapshot.)

\textbf{APM stage:} suppose APM$_1$ (context $c_0$) returns $0.915$ and APM$_2$ (context previous byte $+\ c_0$) returns $0.931$ (hypothetical values). Final probability: $\tfrac14(2\cdot0.9022+0.915+0.931)=0.9126$.

\textbf{Coding:} the bit is $1$: cost $-\log_2 0.9126=0.1319$ bits. The coder narrows its interval by the factor $0.9126$.

\textbf{Learning (with $\eta=0.02$):} the error uses the \emph{mixer's own} output: $y-p_{mix}=0.0978$. Weight changes $\Delta w_i=\eta\cdot0.0978\cdot s_i=(0.0004,\,0.0027,\,0.0051,\,0.0068,\,0.0043,\,0,\,0.0020)$, new weights $(0.1004,\,0.2027,\,0.3051,\,0.2568,\,0.1543,\,0,\,-0.0480)$; with these the same inputs would give $p_{mix}=0.9067$. The counters and StateMaps of each model are also moved toward $y=1$, the APM nodes nearest to the inputs move toward 1, and $c_0\leftarrow 11$.
\end{example}

\section{What dominates the time?}

\begin{itemize}
\item \textbf{Memory latency.} Each model costs about two random cache misses per byte (one per nibble), $\approx100$\,ns each. With 40 models, $\approx8\,\mu$s per byte just waiting for memory unless you \emph{prefetch} (compute next-nibble hashes early and issue prefetch instructions).
\item \textbf{Mixer arithmetic.} $N$ inputs $\times$ $K$ mixers $\times$ (dot product + update) per bit; with $N\approx500$, $K\approx10$ this is $\approx10^4$ multiply-adds per bit, perfectly suited to SIMD (16-bit integer inputs and weights with \code{pmaddwd}-style instructions).
\item \textbf{Neural components.} An LSTM with hidden size $H$ costs $\sim8H(H+I)$ multiply-adds per step for the forward pass alone (Chapter~11), vastly more than everything else if $H$ is in the hundreds.
\end{itemize}

\begin{exercise}[Extend the walk-through]
Suppose the match model became active with expected bit 1 and a learned reliability of $0.9986$, with weight $0.4$. Recompute $z$ and $p_{mix}$. \emph{Answer:} $s=\st(0.9986)=6.570$, $z=2.2219+0.4\cdot6.57=4.850$, $p_{mix}=0.9922$.
\end{exercise}


% =====================================================================
\part{Toward the Hutter Prize}
% =====================================================================

\chapter{Modelling Wikipedia: Text, XML and Preprocessing}

\section{What is in \code{enwik9}?}

\code{enwik9} is the first $10^9$ bytes of the XML dump of the English Wikipedia taken on 3 March 2006. It is not clean prose. Typical content:
\begin{itemize}
\item XML skeleton per article: page, title, id and revision elements, timestamp, contributor, comment, and the text element carrying the article body.
\item MediaWiki markup inside the text: double-bracket links (with optional labels), double-brace templates, bold and italic quote marks, equals-sign headings, tables, lists, categories, redirects.
\item XML/HTML character entities (for quotes, ampersands, angle brackets), because the dump is XML-escaped.
\item English prose, numbers and dates, names, and a lot of non-English words and UTF-8 multibyte sequences.
\end{itemize}
Good compression requires a model for \emph{each} of these ``sub-languages'' and a mixer that knows which one it is currently in. Obtain the file yourself from Mahoney's site (it is about 300\,MB zipped) and \emph{look at it} with a hex viewer and \code{less} before designing anything: a few hours reading the data pays for itself.

\section{Preprocessing: transform the data before modelling}

All top Hutter Prize entries since the early years use a reversible \textbf{preprocessing transform} that rewrites the text into a form that is easier to model; the decompressor undoes it. Typical transforms:
\begin{description}
\item[Capital conversion.] Replace an uppercase letter by a flag byte followed by the lowercase letter (or flag a whole capitalised word). ``The'' and ``the'' now share statistics.
\item[Dictionary (word-replacement) coding.] Replace common words and phrases by short codes (one or two bytes). The text shrinks and the contexts become \emph{more informative per byte}: an order-3 context now spans maybe 3 words instead of 3 letters. This is typically a net win of several percent, and the top entries use large, carefully constructed dictionaries.
\item[XML / markup tokenisation.] Replace tags and common markup by single symbols.
\item[Number and date normalisation, entity decoding, UTF-8 handling.]
\end{description}

\begin{keyidea}[The size of the dictionary counts]
Because $S$ includes the decompressor, a dictionary of $D$ bytes must save more than $D$ bytes of compressed output (and the dictionary itself can be stored compressed). A word that occurs once in \code{enwik9} is never worth putting in the dictionary; a word that occurs a million times certainly is. The optimum is found by experiment, and the transformation is a major tuning area. The preprocessor associated with early Hutter Prize record holders (the ``HP-2017'' preprocessor named in public STARLIT repository documentation) is specific to Wikipedia text.
\end{keyidea}

\begin{example}[A toy dictionary transform]
Original: \texttt{The United States of America was...} Step 1 (capital flag $\uparrow$ = next letter was uppercase): \texttt{$\uparrow$the $\uparrow$united $\uparrow$states of $\uparrow$america was...} Step 2 (dictionary codes $\#n$): \texttt{$\uparrow\#1\ \uparrow\#2\ \uparrow\#3\ \#4\ \uparrow\#5\ \#6$}. The 36-character phrase becomes about 12 symbols, and the model's \emph{order-2} context now covers the equivalent of 2 words. The decoder inverts the steps in reverse order.
\end{example}

\section{Article reordering}

The articles in \code{enwik9} appear in a fixed order unrelated to topic similarity. Compression improves when related articles are \emph{adjacent}, because the match model and the high-order contexts (and the LSTM state) can reuse recent vocabulary and phrasing. The STARLIT submission (2021) exploited exactly this: according to its repository it permutes the articles of \code{enwik9} before compressing with \code{cmix}, and stores the permutation so the decompressor can restore the original order; the cost of storing the permutation is part of $S$. This is a good illustration of the general principle: \textbf{any side information you add must pay for itself}.

\section{Specialised context models for Wikipedia}

On top of the generic orders 1--6 and high orders, text-oriented CM compressors add:
\begin{itemize}
\item \textbf{Word model}: hash of the current word prefix; of the previous word and current prefix; of the two previous words (bigram/trigram); letter-case folded. Capture ``next word given previous words''.
\item \textbf{Sparse and skip contexts}: e.g.\ bytes 2--3 back; useful when a small variation breaks the contiguous context.
\item \textbf{Structure state}: inside a tag name? inside an attribute value? inside \verb|[[ ]]|, \verb|{{ }}|, a table row? depth of nesting. The \emph{state} is used both as a context (hashed with the order-$n$ context) and as the \emph{mixer selector}, so the weights specialise to the structure.
\item \textbf{Link / title echo}: the \verb|<title>| of an article appears again in its first sentence and links; a model can predict the repeated string from the title (a match model restricted to the title is a cheap way).
\item \textbf{Numeric models}: digits following a pattern (years $19xx$, consecutive ids, timestamps incrementing).
\item \textbf{UTF-8 position}: for multibyte sequences, continuation bytes are almost deterministic given the lead byte.
\item \textbf{Indirect contexts and run models} (Chapter 5) at several orders.
\end{itemize}

\begin{handson}[Find where the bits go]
Instrument your compressor to accumulate $-\log_2p$ \emph{per category} of the input: XML tags, markup, English words, numbers, non-ASCII, whitespace. This breakdown tells you where the remaining bytes are. Experienced Hutter Prize participants spend a large fraction of their time on such diagnostics; a model improvement that fixes a category that accounts for 1\% of the output cannot yield more than 1\%.
\end{handson}

\section{Bytes, words or tokens?}

An LLM-style tokenizer (BPE) is, from this point of view, just another dictionary transform; whether BPE-like tokens help a CM depends on whether the gain in context span beats the loss of byte-level sharing between related forms (``walk'', ``walks'', ``walked''). The literature of the prize suggests a hybrid: dictionary-code the very common words, leave the rest as letters.

\section{Tuning on \code{enwik8} versus \code{enwik9}}

\code{enwik8} (the first $10^8$ bytes) is a convenient development set, and \code{enwik9} is the real target. Differences arise: more history makes high-order and match models stronger; the hash tables are under far greater pressure; slowly-adapting parameters (counter limits, learning-rate decay) have a different optimum. Budget time to re-tune on the full file. The prize FAQ explicitly allows (and indeed encourages) tailoring a compressor to \code{enwik9}.

% =====================================================================
\chapter{Neural Components: the LSTM and \code{cmix}}
% =====================================================================

\section{cmix in one paragraph}

\code{cmix} (Byron Knoll) is an open-source compressor in the PAQ tradition that combines a very large number of models (PAQ8-derived context models, many specialised models, and neural networks, notably LSTM byte predictors) into a layered mixing network followed by APM/SSE stages. It is very slow (hundreds of times slower than \code{lpaq}) and achieves among the best published results on text. The recent \code{enwik9} records (\code{fx-cmix}, \code{fx2-cmix}; Table~\ref{tab:records}) are tuned derivatives of the \code{cmix} line. \emph{I describe the architecture in general terms here; you must read the \code{cmix} source on GitHub to learn the actual configuration, which changes between versions.}

\section{The LSTM predictor}

A recurrent network reads the byte stream one symbol at a time and keeps a hidden state that summarises \emph{everything it considers relevant} about the past, including long-range and semantic information that no $n$-gram context can see. The long short-term memory (LSTM) cell, with input $x_t$, previous hidden state $h_{t-1}$ and cell state $c_{t-1}$:
\begin{align*}
f_t&=\sigma(W_fx_t+U_fh_{t-1}+b_f)&&\text{forget gate}\\
i_t&=\sigma(W_ix_t+U_ih_{t-1}+b_i)&&\text{input gate}\\
o_t&=\sigma(W_ox_t+U_oh_{t-1}+b_o)&&\text{output gate}\\
\tilde c_t&=\tanh(W_cx_t+U_ch_{t-1}+b_c)&&\text{candidate}\\
c_t&=f_t\odot c_{t-1}+i_t\odot\tilde c_t&&\text{cell update}\\
h_t&=o_t\odot\tanh(c_t)&&\text{hidden state}
\end{align*}
A softmax output layer turns $h_t$ into a probability distribution $q(\cdot)$ over the 256 possible next bytes.

\paragraph{Size and cost.} With hidden size $H$ and input size $I$ (e.g.\ a one-hot byte, $I=256$), one layer has $4H(H+I+1)$ parameters and one forward step costs about $4H(H+I)$ multiply-adds, plus the output layer $256H$. For $H=256$, $I=256$: $\approx5.3\cdot10^5$ parameters and $\approx5.3\cdot10^5$ multiply-adds ($\approx1$ MFLOP) per byte forward, roughly 2--3 times that to train by back-propagation through time.

\section{From a byte distribution to bit probabilities}

The arithmetic coder consumes \emph{bit} probabilities, so a byte-level predictor must be converted. For the bit at tree node $c_0$, the probability that it is 1 is the mass of byte values consistent with the prefix \emph{and} the bit 1, divided by the mass consistent with the prefix:
\[
P(y_k=1\mid\text{prefix})=\frac{\sum_{b\in B_1(\text{prefix})}q(b)}{\sum_{b\in B(\text{prefix})}q(b)}.
\]

\begin{example}[Two-bit toy alphabet]
Alphabet $\{00,01,10,11\}$ with network output $q=(0.5,0.2,0.2,0.1)$.
\begin{itemize}
\item First bit: $P(1)=q(10)+q(11)=0.3$.
\item Second bit given first bit $0$: $P(1\mid0)=\dfrac{q(01)}{q(00)+q(01)}=\dfrac{0.2}{0.7}=0.2857$.
\item Second bit given first bit $1$: $P(1\mid1)=\dfrac{q(11)}{q(10)+q(11)}=\dfrac{0.1}{0.3}=0.3333$.
\end{itemize}
Check: the probability of the symbol $10$ is $P(\text{first}=1)\cdot P(\text{second}=0\mid1)=0.3\cdot(1-0.3333)=0.2=q(10)$. The decomposition reproduces the original distribution exactly.
\end{example}

These bit probabilities are then \emph{stretched and fed to the mixer like any other model's prediction}; the mixer learns how much to trust the LSTM, bit by bit and in each selector context. This hybrid is the essential design: the LSTM is an additional expert, complementary to the count-based models.

\section{Online training}

There is no pre-training. The network is trained \emph{while compressing}: after each byte (or every few bytes) the actual byte is the training target and the weights are updated by truncated back-propagation through time over a window of recent bytes, using an optimiser such as Adam or RMSprop-style normalised steps (as far as I know \code{cmix} uses an Adam-like optimiser; check the source). The decompressor does the identical computation in the identical order, so it re-learns exactly the same weights.

\begin{keyidea}[Learned parameters are free; computation is not]
Because the LSTM weights are learned on the fly from a fixed initialisation (a seed), they are \emph{not part of $S$}. A pre-trained network of the same size would have to be shipped and would cost its full size in bytes. The Prize therefore rewards models that learn quickly from a single pass over 1\,GB, within the time limit. This is why the model class that wins is ``online-learning ensemble'' rather than ``big pretrained model'', and why a general-purpose LLM, however good at English, does not win (the prize FAQ mentions the compute limits; recent work such as \emph{Language Modeling Is Compression} (Del\'etang et al., 2023) makes the same point about model size counting).
\end{keyidea}

\section{Budgeting compute for a neural model}

\begin{example}[How much LSTM fits?]
Allowed: 50 hours for $10^9$ bytes $=180\ \mu$s/byte. On one core, suppose a sustained $10$--$30$\,GFLOP/s with well-vectorised code (single-core AVX2 peak is higher but memory bound code is slower). Then the \emph{whole} model gets $1.8$--$5.4$ MFLOP per byte. A single-layer LSTM with $H=256$ needs about $1$ MFLOP forward and $2$--$3$ MFLOP backward per byte: it would consume most of the budget by itself, leaving little for the 100+ context models. Hence the real systems use cleverness: smaller hidden sizes, training only every few bytes, sparse or low-precision updates, sharing the cost across the model's inputs. Estimating this budget before writing any code is a critical design step.
\end{example}

\section{Other neural ideas}

\begin{itemize}
\item \textbf{Bigger or deeper recurrent models} (multi-layer LSTMs, GRUs) and \textbf{Transformers} trained online. Fabrice Bellard's NNCP project is, as I recall, the best-known system doing lossless compression with Transformer/LSTM networks and ranks near the top of Mahoney's Large Text Compression Benchmark, but it does not fit the Prize's time limits; check its documentation for current numbers.
\item \textbf{Neural mixers.} Replace the logistic mixer with a small neural network that also sees context features.
\item \textbf{Mixing at the byte level}: combine byte-level distributions (from the LSTM) before decomposing to bits.
\item \textbf{Distilling} expensive models into cheap ones is not allowed to cheat the limits, but \emph{architecture search} on a development subset is.
\end{itemize}

\section{Floating-point determinism}

Compression and decompression must perform \emph{bit-identical} computations. With integer-only CM this is automatic; with floating-point neural networks you must ensure: the same binary and the same CPU instruction paths on both sides; a fixed order of summation (no multithreaded reductions with non-deterministic order); no dependence on uninitialised memory or on timing; and care with \code{-ffast-math} and library \code{exp}/\code{tanh} implementations. The safest course is to use your own implementations of \code{exp}, \code{tanh} and \code{sigmoid} (polynomial or table based) in fixed code.

\begin{exercise}[Byte-to-bit decomposition]
A network outputs $q$ over 8 byte values (3-bit symbols $000..111$): $(0.3,0.1,0.1,0.1,0.2,0.1,0.05,0.05)$. Compute the three bit probabilities for coding the symbol $101$. \emph{Answer:} first bit: $P(1)=0.2+0.1+0.05+0.05=0.4$; given prefix $1$, second bit: $P(1\mid1)=(0.05+0.05)/0.4=0.25$; given prefix $10$, third bit: $P(1\mid10)=0.1/(0.2+0.1)=0.3333$. The coded bits are $1,0,1$, so the symbol's probability is $0.4\cdot(1-0.25)\cdot0.3333=0.1=q(101)$, as it must be.
\end{exercise}


% =====================================================================
\chapter{Engineering Under the Prize Limits}
% =====================================================================

\section{The constraints, restated as design requirements}

\begin{center}
\begin{tabular}{@{}ll@{}}
\toprule
constraint (check the current rules!) & design consequence\\
\midrule
$\approx50$ h, single core & $\le180\,\mu$s per byte for the whole model; no multithreading\\
$\approx10$ GB RAM & table sizes, bucket formats, LSTM state, history buffer\\
$S=$ decompressor + archive & code size and dictionary size count\\
decoder $=$ encoder cost & you pay the time twice when testing\\
exact reproduction on their machine & determinism; no machine-specific behaviour\\
\bottomrule
\end{tabular}
\end{center}

\section{Speed engineering}

\begin{itemize}
\item \textbf{Cache discipline.} Place all the states for one nibble of one context in a single 64-byte cache line. Compute the hashes for the next nibble as soon as the current one is known, and issue \code{\_\_builtin\_prefetch}.
\item \textbf{SIMD for mixers.} Keep inputs and weights as 16-bit integers; use multiply-add instructions that accumulate pairs into 32 bits (e.g.\ AVX2 \code{vpmaddwd}). Align arrays to 32 bytes. Mixer cost then becomes a few cycles per 16 inputs.
\item \textbf{Avoid division and \code{exp}} in the inner loop: use lookup tables for stretch and squash.
\item \textbf{Profile with real counters} (\code{perf}): memory stalls and branch mispredictions rather than instruction count are usually what limits CM.
\item \textbf{Two-pass tricks}: the encoder may analyse the file (choose a dictionary, a parameter set, an article order) and transmit the result; the decoder just uses it. Analysis time on the encoder side is not repeated by the decoder, but the total encoder limit still applies.
\end{itemize}

\section{Memory engineering}

\begin{itemize}
\item Count bytes per context slot honestly and give each model a budget; the best split is found empirically.
\item Use \emph{checked} buckets with sensible replacement (Chapter 4).
\item Huge pages (e.g.\ 2\,MB) reduce TLB misses in multi-gigabyte tables.
\item Bound the history buffer; for \code{enwik9} the whole file ($10^9$ bytes) can stay in memory for the match model.
\end{itemize}

\section{Correctness engineering}

A compressor that cannot decompress is worth nothing, and bugs in CM are subtle because the encoder and decoder share almost all code. A disciplined test regime:
\begin{enumerate}
\item A \textbf{round-trip test on every change}, on a small file (1 MB) and a medium file (100 MB). Compare with \code{cmp} or \code{sha1sum}.
\item \textbf{Determinism test}: run compression twice and require byte-identical archives. Run decompression on a different (but same-architecture) machine.
\item \textbf{Never let encoder-only information leak into the model}: for example, using the \emph{next} byte or the file length while predicting. The decoder cannot do that. A classic way to catch this: structure the code so that a single \code{predict()/update(bit)} interface is used by both encoder and decoder.
\item \textbf{Edge cases}: empty input, 1-byte input, runs of identical bytes, all 256 byte values, files smaller than the hash window.
\item \textbf{Sanitisers} (\code{-fsanitize=address,undefined}) on the small test; uninitialised memory is a common source of non-determinism.
\end{enumerate}

\begin{pitfall}[Floating point and compilers]
Changing the compiler version or flags can change the results of floating-point code (fused multiply-add contraction, vectorised reductions). If both compressor and decompressor are built from the same source with the same flags on the same machine they will agree; but any rebuild between compress and decompress runs risks catastrophic failure. For the Prize, the decompressor is shipped as a self-extracting executable, which removes the risk at submission time, but \emph{you} must check by running that exact executable.
\end{pitfall}

\section{Development workflow}

\begin{enumerate}
\item \textbf{Fast loop:} first 10--20 MB of \code{enwik8} (seconds to minutes). Use for algorithmic ideas; ignore changes below $0.1\%$ unless cheap.
\item \textbf{Medium loop:} full \code{enwik8} or the first 100--200 MB of \code{enwik9} (hours); run many experiments in parallel on a multi-core development machine. (The \emph{submission} is single-core but your \emph{development} need not be.)
\item \textbf{Full loop:} complete \code{enwik9}, 50 hours; do this rarely and always with a final-candidate build, and keep logs (size, time, peak memory) for each.
\item Log everything: git commit hash, parameters, size to the byte, time, memory.
\end{enumerate}

\section{Parameter tuning}

A CM compressor has hundreds of numeric parameters (counter limits, learning rates, table sizes, selector definitions). Practical approaches:
\begin{itemize}
\item \textbf{Coordinate descent}: change one parameter at a time up and down by a factor, keep improvements. Slow but robust.
\item \textbf{Random search / evolutionary} on the medium loop with many cores. The record holders reportedly relied heavily on automated tuning.
\item \textbf{Bayesian optimisation} for small numbers of expensive parameters.
\end{itemize}
Beware \emph{noise from interactions}: parameters tuned on 20 MB may be wrong at 1 GB (Chapter~10); re-verify the big ones on the medium loop.

\section{Shrinking the decompressor}

The decompressor binary and its data (dictionaries, initial tables) count. Techniques: compile with \code{-Os} versus \code{-O3} experimentally (speed versus size); strip symbols; compress the dictionary with the model itself (bootstrapping: decode the dictionary using a smaller model first); remove unused code paths (the decoder never needs the encoder's analysis code).

\begin{exercise}[Budget your cache misses]
With 80 context models, each needing 2 random DRAM accesses per byte at 90\,ns without prefetching, compute the memory time per byte and compare with the 180\,$\mu$s budget. \emph{Answer:} $80\cdot2\cdot90\text{ ns}=14.4\,\mu$s ($8\%$ of the budget); with prefetching most of this overlaps with computation. Memory latency is \emph{not} the main problem at this model count; compute (mixers, neural networks) is.
\end{exercise}

% =====================================================================
\chapter{Designing Your Own Algorithm}
% =====================================================================

\section{An honest assessment}

To receive any money you must beat the current record by at least $1\%$ ($\approx1.1$ MB) and the record has been advanced mostly by large, long-developed systems and by a few people who understand every line of them. It is a legitimate and valuable goal; a realistic plan is staged, with each stage useful in itself (a lecture, a paper, a benchmark entry on Mahoney's \emph{Large Text Compression Benchmark}) even if the final prize is out of reach.

\section{Staged plan}

\begin{enumerate}[label=\textbf{Stage \arabic*.},leftmargin=*]
\item \textbf{Reproduce.} Understand and run the code of Appendix~A; reproduce the ablation table; implement the same ideas yourself from a blank file. Read \code{lpaq1} completely (about 500 lines).
\item \textbf{PAQ8 level.} Add: bit-history state machines and StateMaps, nibble-bucketed checked hash tables, two-layer mixers, several APMs, run/indirect/sparse/word models, several match models. Validate against \code{lpaq}/\code{paq8} binaries on \code{enwik8}.
\item \textbf{Text specialisation.} XML/markup model, capital conversion, dictionary preprocessor, article reordering experiments. Build the per-category bit-cost diagnostics.
\item \textbf{Neural component.} Add an LSTM (start small), feeding the mixer. Measure gain per millisecond.
\item \textbf{Match the state of the art.} Run \code{cmix} (and, if you can, its \code{fx} variants) on the same data. Until you match them on your development set you have nothing new to contribute.
\item \textbf{Engineer for the limits.} Profile, vectorise, prefetch; confirm the 50-hour, single-core, 10\,GB envelope on a clean machine.
\item \textbf{Innovate.} Only now explore genuinely new ideas, ranked by (expected gain)/(effort).
\end{enumerate}

\section{A checklist for evaluating every idea}

For each proposed change answer, \emph{in writing}:
\begin{enumerate}
\item Which category of the bit-cost breakdown does it target, and how many bytes of output does that category currently account for? (This is the upper bound on the gain.)
\item What does it cost in time per byte and in memory?
\item What does it cost in decompressor size?
\item Does the gain survive at the large scale (medium loop)?
\item Is the gain \emph{additive} with the others? (Run the ablation with and without.)
\end{enumerate}

\section{Where gains may still be hiding}

These are research directions, not promises; every item has probably been tried in some form:
\begin{itemize}
\item Better \textbf{tokenisation/dictionaries} jointly optimised with the model; learned dictionaries; handling of rare words and names.
\item \textbf{Better counters/state machines} for \code{enwik}-style data; two-rate counters; state machines trained from data rather than hand-designed.
\item \textbf{Mixer architecture}: richer selectors (grammatical state, title match), three layers, neural mixers, better learning-rate schedules.
\item \textbf{SSE} contexts derived from the LSTM state or the match state.
\item \textbf{Neural models}: bigger or better-structured recurrent models; attention over the recent history; mixture of small networks specialised by structure.
\item \textbf{Two-pass modelling}: encoder-side search over parameters and orderings, transmitting a few kilobytes of side information.
\item \textbf{Precision}: higher-resolution probabilities in the final stages and the coder (Chapter 3).
\item \textbf{Exploiting the Wikipedia 2006 specifics}: templates, categories and interlanguage links have regularities that dedicated models can predict.
\end{itemize}

\section{Social and practical advice}

\begin{itemize}
\item Join the discussion forum linked from the prize page, and read the threads of previous winners: they often describe the changes that mattered.
\item Read the licences: the Prize requires an OSI-approved open-source licence for the winning submission, and building on \code{cmix}/\code{paq8} code means complying with their licences.
\item Re-read the official rules before you begin, before you spend 50 hours on a run, and before you submit. They determine eligibility, hardware and the 30-day comment period.
\end{itemize}

\begin{exercise}[Where can the gain come from?]
Suppose your breakdown shows XML tags use $2\%$ of the compressed output, markup syntax $8\%$, English words $70\%$, numbers $6\%$, other $14\%$. To gain $1\%$ overall, by what percentage must you reduce the English-words category? the XML-tags category? \emph{Answer:} words: $1/70=1.43\%$; tags: $1/2=50\%$ (implausible). Prioritise the large categories unless the small one is cheap.
\end{exercise}


% =====================================================================
\appendix
% =====================================================================

\chapter{A Complete Teaching Compressor (tinycm)}
\label{app:code}

This is a complete, working, readable context-mixing compressor, about 190 lines of C. It is \emph{not} competitive with \code{paq8} or \code{cmix}; it exists so that every idea in Part~II can be seen in running code and modified. I compiled it with \code{gcc -O2} and verified that decompression reproduces the input exactly (byte-for-byte \code{cmp}).

\section{Architecture}

\begin{center}
\begin{tabular}{@{}lp{10.5cm}@{}}
\toprule
component & implementation\\
\midrule
context models & orders 1, 2, 3, 4, 6 and a word model, each hashed into a $2^{22}$-slot table; order 0 direct\\
slot & 32 bits: 22-bit probability and 10-bit count, adaptive rate $1/(n+1.5)$, limit 255\\
match model & hash of last 6 bytes, longest-recent-match pointer; reliability learned per length bucket\\
mixer & single layer, 9 inputs ($\st$ of 8 models and a bias), 256 weight sets selected by $c_0$\\
APM stages & two: context $c_0$ and (previous byte, $c_0$); averaged $2{:}1{:}1$ with the mixer output\\
coder & 32-bit carry-less binary arithmetic coder, 12-bit probabilities\\
\bottomrule
\end{tabular}
\end{center}

Omitted on purpose (they are exercises): bit-history states and StateMaps (the slots hold direct probabilities), checksummed buckets, multiple mixers and a second layer, run/indirect/sparse models.

\section{Results on a 1\,MB sample}

\begin{table}[h]
\centering
\begin{tabular}{@{}lrr@{}}
\toprule
compressor & size of 1\,000\,000-byte sample & bits per byte\\
\midrule
gzip -9 & 214\,589 & 1.717\\
bzip2 -9 & 134\,193 & 1.074\\
xz -9e & 102\,576 & 0.821\\
\textbf{tinycm} (this appendix) & \textbf{89\,778} & \textbf{0.718}\\
\bottomrule
\end{tabular}
\caption{The sample is 1\,MB of English documentation and licence text taken from the preparation machine (not Wikipedia; highly repetitive). tinycm runs at about 1.4 s per MB for both compression and decompression ($\approx0.7$ MB/s).}
\end{table}

For orientation: at this speed $10^9$ bytes would take about 20 minutes; the Prize allows 50 hours, a factor of 150 more.

\section{Source code}

\begin{lstlisting}[style=ccode]
/* tinycm.c -- a small, readable context-mixing compressor (teaching code).
 * Usage:  tinycm c in out   |   tinycm d in out                         */
#include <stdio.h>
#include <stdlib.h>
#include <string.h>
#include <stdint.h>
#include <ctype.h>
typedef uint8_t U8; typedef uint32_t U32; typedef uint64_t U64; typedef uint16_t U16;

/* ---------- logistic helpers: squash(x)=4096/(1+e^-x/256), stretch=inverse ---------- */
static int squash(int d) {
  static const int t[33] = {1,2,3,6,10,16,27,45,73,120,194,310,488,747,1101,
    1546,2047,2549,2994,3348,3607,3785,3901,3975,4022,4050,4068,4079,4085,
    4089,4092,4093,4094};
  if (d > 2047) return 4095;
  if (d < -2047) return 0;
  int w = d & 127; d = (d >> 7) + 16;
  return (t[d] * (128 - w) + t[d + 1] * w + 64) >> 7;
}
static short stretch_t[4096];
static void init_stretch(void) {
  int pi = 0;
  for (int x = -2047; x <= 2047; x++) {
    int v = squash(x);
    for (int i = pi; i <= v; i++) stretch_t[i] = x;
    pi = v + 1;
  }
  for (int i = pi; i < 4096; i++) stretch_t[i] = 2047;
}
#define stretch(p) stretch_t[p]

/* ---------- adaptive probability slot: 22-bit probability + 10-bit count ---------- */
static int dt[1024];
static inline int  slot_p(U32 s) { return s >> 20; }            /* 12-bit probability */
static inline void slot_update(U32 *s, int y, int limit) {
  int n = *s & 1023, p = *s >> 10;
  if (n < limit) ++*s; else *s = (*s & 0xfffffc00) | limit;
  int64_t delta = (int64_t)(((y << 22) - p) >> 3) * dt[n];
  *s += (U32)delta & 0xfffffc00;
}

/* ---------- APM / SSE stage: 33 interpolation nodes per context ---------- */
typedef struct { U16 *t; int idx; } APM;
static void apm_init(APM *a, int n) {
  a->t = malloc(n * 33 * sizeof(U16));
  for (int i = 0; i < n; i++)
    for (int j = 0; j < 33; j++) a->t[i * 33 + j] = squash((j - 16) * 128) * 16;
}
static int apm_pp(APM *a, int pr, int cx) {
  int s = (stretch(pr) + 2048) * 32, lo = s >> 12, wt = s & 4095;
  a->idx = cx * 33 + lo + (wt >> 11);
  return (a->t[cx * 33 + lo] * (4096 - wt) + a->t[cx * 33 + lo + 1] * wt) >> 16;
}
static void apm_update(APM *a, int y, int rate) {
  int g = (y << 16) + (y << rate) - y - y;
  a->t[a->idx] += (g - a->t[a->idx]) >> rate;
}

/* ---------- the predictor ---------- */
#define NCM 6                 /* hashed context models: o1 o2 o3 o4 o6 word  */
#define NIN (NCM + 3)         /* + order-0, match model, bias                */
#define TBITS 22
#define HBITS 22              /* match-model hash table size                 */
#define BUFSZ (1 << 26)
static U32 *T[NCM], o0[256], mm_sm[64 * 2];
static U32 ctx[NCM], idx[NCM], *o0p, *mmp;
static int st[NIN], wsel;
static int *W;                /* weight sets: 256 sets x NIN weights (16.16)  */
static APM a1, a2;
static int c0 = 1, bitpos = 0, pr = 2048, pr_mix = 2048;
static U64 c8 = 0;            /* last 8 bytes                                */
static U32 wordh = 0;
static U8 *buf; static U32 pos = 0, mptr = 0, mlen = 0, *mht;
static int mexp = -1;         /* expected bit from the match model, or -1    */

static void predictor_init(void) {
  init_stretch();
  for (int i = 0; i < 1024; i++) dt[i] = 16384 / (i + i + 3);
  for (int i = 0; i < NCM; i++) {
    T[i] = malloc((size_t)4 << TBITS);
    for (U32 j = 0; j < (1u << TBITS); j++) T[i][j] = 1u << 31;
  }
  for (int i = 0; i < 256; i++) o0[i] = 1u << 31;
  for (int i = 0; i < 128; i++) mm_sm[i] = 1u << 31;
  W = malloc(256 * NIN * sizeof(int));
  for (int i = 0; i < 256 * NIN; i++) W[i] = 1 << 14;      /* 0.25 */
  apm_init(&a1, 256); apm_init(&a2, 65536);
  buf = calloc(BUFSZ, 1); mht = calloc(1u << HBITS, 4);
}
static inline U32 hidx(U32 h, U32 c) {
  U32 x = (h + c * 0x9E3779B1u) * 0x85EBCA6Bu;
  x ^= x >> 15; x *= 0xC2B2AE35u;
  return x >> (32 - TBITS);
}
static int predict(void) {
  for (int i = 0; i < NCM; i++) {
    idx[i] = hidx(ctx[i], c0);
    st[i] = stretch(slot_p(T[i][idx[i]]));
  }
  o0p = &o0[c0]; st[NCM] = stretch(slot_p(*o0p));
  /* match model: only active if the bits seen so far agree with prediction */
  mexp = -1;
  if (mlen > 0 && ((buf[mptr] | 256) >> (8 - bitpos)) == c0) {
    mexp = (buf[mptr] >> (7 - bitpos)) & 1;
    int l = mlen > 31 ? 31 : mlen;
    mmp = &mm_sm[l * 2 + mexp];
    st[NCM + 1] = stretch(slot_p(*mmp));
  } else { mlen = 0; st[NCM + 1] = 0; }
  st[NCM + 2] = 256;                                         /* bias */
  wsel = c0 * NIN;
  int64_t dot = 0;
  for (int i = 0; i < NIN; i++) dot += (int64_t)st[i] * W[wsel + i];
  pr_mix = squash((int)(dot >> 16));
  int p1 = apm_pp(&a1, pr_mix, c0);
  int p2 = apm_pp(&a2, pr_mix, (c0 | (int)(c8 & 255) << 8) & 0xffff);
  pr = (2 * pr_mix + p1 + p2 + 2) >> 2;
  if (pr < 1) pr = 1;
  if (pr > 4095) pr = 4095;
  return pr;
}
static void update(int y) {
  for (int i = 0; i < NCM; i++) slot_update(&T[i][idx[i]], y, 255);
  slot_update(o0p, y, 255);
  if (mexp >= 0) slot_update(mmp, y, 1023);
  int err = ((y << 12) - pr_mix) * 6;                         /* mixer learning rate */
  for (int i = 0; i < NIN; i++) W[wsel + i] += (st[i] * err + 0x8000) >> 16;
  apm_update(&a1, y, 7); apm_update(&a2, y, 7);
  c0 = c0 * 2 + y; bitpos++;
  if (mexp >= 0 && mexp != y) mlen = 0;
  if (bitpos == 8) {                                          /* byte finished */
    int c = c0 & 255;
    buf[pos++ & (BUFSZ - 1)] = c;
    c8 = (c8 << 8) | c;
    wordh = isalpha(c) ? (wordh + tolower(c) + 1) * 0x2F0B4A27u : 0;
    static const U64 mask[5] = {0xff, 0xffff, 0xffffff, 0xffffffffULL, 0xffffffffffffULL};
    for (int i = 0; i < 5; i++) ctx[i] = (U32)(((c8 & mask[i]) * 0x9E3779B97F4A7C15ULL) >> 32) + i;
    ctx[5] = wordh * 0x3D4D51CBu + 99;
    /* match model: follow current match or look up the last 6 bytes */
    if (mlen > 0 && buf[mptr] == c) { mlen++; mptr++; } else mlen = 0;
    U32 h = (U32)(((c8 & 0xffffffffffffULL) * 0x9E3779B97F4A7C15ULL) >> (64 - HBITS));
    if (mlen == 0) {
      mptr = mht[h];
      if (mptr > 0) while (mlen < 255 && mlen < mptr &&
             buf[(mptr - mlen - 1) & (BUFSZ-1)] == buf[(pos - mlen - 1) & (BUFSZ-1)]) mlen++;
    }
    mht[h] = pos;
    c0 = 1; bitpos = 0;
  }
}

/* ---------- binary arithmetic coder (carry-less, 32-bit) ---------- */
static U32 x1 = 0, x2 = 0xffffffff, x = 0;
static FILE *fin, *fout;
static void encode(int y, int p) {
  U32 xmid = x1 + ((x2 - x1) >> 12) * p;
  y ? (x2 = xmid) : (x1 = xmid + 1);
  while (((x1 ^ x2) & 0xff000000) == 0) { putc(x2 >> 24, fout); x1 <<= 8; x2 = (x2 << 8) | 255; }
}
static int decode(int p) {
  U32 xmid = x1 + ((x2 - x1) >> 12) * p;
  int y = x <= xmid;
  y ? (x2 = xmid) : (x1 = xmid + 1);
  while (((x1 ^ x2) & 0xff000000) == 0) {
    x1 <<= 8; x2 = (x2 << 8) | 255; x = (x << 8) | (getc(fin) & 255);
  }
  return y;
}
int main(int argc, char **argv) {
  if (argc != 4) { fprintf(stderr, "usage: tinycm c|d in out\n"); return 1; }
  fin = fopen(argv[2], "rb"); fout = fopen(argv[3], "wb");
  if (!fin || !fout) { perror("open"); return 1; }
  predictor_init();
  if (argv[1][0] == 'c') {
    fseek(fin, 0, SEEK_END); U32 n = ftell(fin); fseek(fin, 0, SEEK_SET);
    fwrite(&n, 4, 1, fout);
    for (U32 i = 0; i < n; i++) {
      int c = getc(fin);
      for (int b = 7; b >= 0; b--) { int y = (c >> b) & 1; encode(y, predict()); update(y); }
    }
    for (int i = 0; i < 4; i++) { putc(x1 >> 24, fout); x1 <<= 8; }  /* flush */
  } else {
    U32 n; fread(&n, 4, 1, fin);
    for (int i = 0; i < 4; i++) x = (x << 8) | (getc(fin) & 255);
    for (U32 i = 0; i < n; i++) {
      int c = 0;
      for (int b = 0; b < 8; b++) { int y = decode(predict()); update(y); c = c * 2 + y; }
      putc(c, fout);
    }
  }
  fclose(fout); return 0;
}

\end{lstlisting}

\section{Suggested modifications (a research-training programme)}

\begin{enumerate}
\item \textbf{Measure.} Reproduce the ablation table of Chapter~6 (Table~6.2) on your own data. Add a per-component bit-cost log.
\item Replace slots by \textbf{bit-history states + StateMap} (Chapter~5) and measure.
\item Add \textbf{checksummed buckets} and test with a file larger than the table can hold.
\item Add \textbf{two more mixers} with selectors (match length; previous byte) and a \textbf{final mixer} (Chapter~6.8).
\item Add a second input per model, $(p-\tfrac12)$, and confidence inputs.
\item Add a \textbf{sparse} model and a \textbf{second match model} with minimum length 12.
\item Replace the hash $\to$ slot design by one slot per \emph{nibble} (Chapter~4.5).
\item Add a small \textbf{LSTM} in floating point and feed its bit-probability to the mixer.
\item Make mixer learning rates \textbf{decay}; add RMS-normalised updates.
\item Write the \textbf{pre-processing} transform of Chapter~10 and measure on a Wikipedia sample.
\end{enumerate}

\chapter{Notation}

\begin{center}
\begin{tabular}{@{}lp{11cm}@{}}
\toprule
symbol & meaning\\
\midrule
$y,\ y_t$ & the bit being coded (0 or 1)\\
$c_0$ & partial-byte context: leading 1 followed by the bits of the current byte seen so far\\
$p_i$ & probability that $y=1$ predicted by model $i$\\
$\st(p)=\ln\frac p{1-p}$ & logit; in the code scaled by 256 and clamped to $\pm2047$\\
$\sq(x)=1/(1+e^{-x})$ & logistic function (inverse of stretch)\\
$s_i=\st(p_i)$ & mixer input\\
$w_{m,i}$ & mixer weight for input $i$ in weight set $m$\\
$\eta$ & mixer learning rate\\
$\ell$ & code length, $-\log_2$ of the probability assigned to what happened\\
$H,\ \KL$ & entropy, Kullback--Leibler divergence\\
$S,\ L,\ Z$ & Prize total size, previous record, prize fund\\
\bottomrule
\end{tabular}
\end{center}

\chapter{Glossary}

\begin{description}
\item[APM / SSE] adaptive probability map / secondary symbol estimation: a learned calibration of a probability, conditioned on a small context.
\item[Arithmetic coding] entropy coding that turns a sequence of probabilities into a number inside nested intervals.
\item[Bit history / state] a one-byte summary of the sequence of bits seen in a context.
\item[Context] any function of the past used to index statistics.
\item[Context mixing (CM)] combining the predictions of many context models with an adaptively trained function.
\item[Indirect model] a model whose context is the history of what followed an earlier context, or which maps a bit history to a probability through a learned table.
\item[Logistic mixing] $p=\sq(\sum w_is_i)$; a product of experts.
\item[Match model] predicts the next symbol by following the most recent earlier occurrence of the current long context.
\item[Online learning] parameters updated while the data stream by; no separate training phase.
\item[StateMap] table mapping a bit-history state (or context) to an adaptively learned probability.
\item[Stretch / squash] logit and logistic functions.
\item[Weight set] one of many weight vectors of a mixer, selected by a context.
\end{description}

\chapter{Suggested Reading and Code to Study}

\paragraph{Read the source code in this order.}
\begin{enumerate}
\item \code{lpaq1} (Matt Mahoney): about 500 lines, the cleanest complete CM compressor; compare every function with this document.
\item \code{paq8} family (\code{paq8l}, and the actively developed \code{paq8px}): the full model zoo.
\item \code{cmix} (Byron Knoll): the PAQ8 models plus LSTM and the layered mixer.
\item The \code{starlit} repository (Margaritov): the article-reordering approach and how the prize constraints were met.
\end{enumerate}

\begin{thebibliography}{99}
\bibitem{shannon} C.~E.~Shannon, ``A mathematical theory of communication'', \emph{Bell System Technical Journal} 27 (1948).
\bibitem{wnc} I.~H.~Witten, R.~M.~Neal, J.~G.~Cleary, ``Arithmetic coding for data compression'', \emph{Communications of the ACM} 30(6), 1987.
\bibitem{ppm} J.~G.~Cleary and I.~H.~Witten, ``Data compression using adaptive coding and partial string matching'', \emph{IEEE Trans.\ Communications} 32(4), 1984.
\bibitem{cover} T.~M.~Cover and J.~A.~Thomas, \emph{Elements of Information Theory}, Wiley.
\bibitem{dce} M.~Mahoney, \emph{Data Compression Explained}, online book, \url{http://mattmahoney.net/dc/dce.html}.
\bibitem{adaptive} M.~Mahoney, ``Adaptive weighing of context models for lossless data compression'', Florida Tech.\ technical report CS-2005-16.
\bibitem{ltcb} M.~Mahoney, \emph{Large Text Compression Benchmark}, \url{http://mattmahoney.net/dc/text.html}.
\bibitem{prize} M.~Hutter et al., \emph{500\,000\texteuro{} Prize for Compressing Human Knowledge} (rules, FAQ, record table), \url{http://prize.hutter1.net/}.
\bibitem{uai} M.~Hutter, \emph{Universal Artificial Intelligence}, Springer, 2005.
\bibitem{poe} G.~E.~Hinton, ``Training products of experts by minimizing contrastive divergence'', \emph{Neural Computation} 14(8), 2002.
\bibitem{lstm} S.~Hochreiter and J.~Schmidhuber, ``Long short-term memory'', \emph{Neural Computation} 9(8), 1997.
\bibitem{adam} D.~P.~Kingma and J.~Ba, ``Adam: a method for stochastic optimization'', ICLR 2015.
\bibitem{ans} J.~Duda, ``Asymmetric numeral systems: entropy coding combining speed of Huffman coding with compression rate of arithmetic coding'', arXiv:1311.2540, 2013.
\bibitem{lmc} G.~Del\'etang et al., ``Language modeling is compression'', arXiv:2309.10668 (ICLR 2024).
\bibitem{cmix} B.~Knoll, \emph{cmix}, \url{https://github.com/byronknoll/cmix}.
\bibitem{nncp} F.~Bellard, \emph{NNCP: lossless data compression with neural networks}, \url{https://bellard.org/nncp/}.
\bibitem{starlit} A.~Margaritov, \emph{STARLIT}, \url{https://github.com/amargaritov/starlit}.
\end{thebibliography}

\end{document}

